\documentclass[aspectratio=169,11pt]{beamer}

% =========================================================
% CSE 219: Signals and Linear Systems
% Introductory slides
% Compile with: pdflatex cse219_intro_slides.tex
% =========================================================

\usepackage[T1]{fontenc}
\usepackage[utf8]{inputenc}
\usepackage{lmodern}
\usepackage{graphicx}
\usepackage{amsmath,amssymb,bm}
\usepackage{booktabs}
\usepackage{tikz}
\usepackage[most]{tcolorbox}
\usepackage{pgfplots}
\usepackage{xcolor}
\usepackage{array}
\usepackage{fontawesome5}
\usepackage{microtype}
\usepackage{multirow}
\usepackage{hyperref}
\usepackage{ragged2e}

\pgfplotsset{compat=1.18}
\usetikzlibrary{arrows.meta,positioning,fit,shapes.geometric,calc,decorations.pathreplacing}

\usepackage{animate}




% =========================================================
% SIGNAL-TRANSFORMATION DRAWING HELPERS
% Put this in the PREAMBLE, before \begin{document}.
% =========================================================

\newcommand{\IrregularCTPath}{%
  (-1.80,0.15)
  .. controls (-1.55,0.70) and (-1.18,0.82) .. (-0.82,0.58)
  .. controls (-0.58,0.38) and (-0.47,-0.38) .. (-0.08,-0.32)
  .. controls (0.12,-0.28) and (0.22,0.78) .. (0.55,0.95)
  .. controls (0.86,0.60) and (1.08,0.20) .. (1.55,0.18)
  .. controls (1.78,0.18) and (1.90,-0.16) .. (2.02,-0.03)
}

\newcommand{\DrawDTSequence}[2]{%
  \foreach \k/\v in {
    -3/0.30,
    -2/0.85,
    -1/0.45,
     0/-0.65,
     1/-0.20,
     2/0.70,
     3/0.30
  }{%
    \draw[very thick,#2] ({\k+#1},0) -- ({\k+#1},\v);
    \fill[#2] ({\k+#1},\v) circle (2.1pt);
    \fill[black] ({\k+#1},0) circle (0.9pt);
  }%
}



% Dense discrete-time sequence:
% #1 = integer shift, #2 = stem color
\newcommand{\DrawDenseDTSequence}[2]{%
  \foreach \k/\v in {
    -7/0.18,
    -6/0.36,
    -5/0.58,
    -4/0.84,
    -3/0.67,
    -2/0.28,
    -1/-0.18,
     0/-0.62,
     1/-0.48,
     2/-0.12,
     3/0.35,
     4/0.72,
     5/0.89,
     6/0.70,
     7/0.42
  }{%
    \draw[very thick,#2] ({\k+#1},0) -- ({\k+#1},\v);
    \fill[#2] ({\k+#1},\v) circle (1.55pt);
    \fill[black] ({\k+#1},0) circle (0.75pt);
  }%
}









% =========================================================
% DESIGN CONTROL PANEL
% =========================================================

% -------------------- color palette --------------------
\definecolor{CourseBlue}{HTML}{2563EB}
\definecolor{CourseNavy}{HTML}{0F172A}
\definecolor{CourseSky}{HTML}{EAF2FF}
\definecolor{CourseGreen}{HTML}{059669}
\definecolor{CourseMint}{HTML}{E8FFF5}
\definecolor{CourseOrange}{HTML}{EA580C}
\definecolor{CourseMaroon}{HTML}{7A0000}
\definecolor{CourseRed}{HTML}{DC2626}

\definecolor{CourseGray}{HTML}{475569}
\definecolor{CourseLight}{HTML}{F8FAFC}
\definecolor{CourseLine}{HTML}{CBD5E1}
\definecolor{CourseSoft}{HTML}{F1F5F9}

\colorlet{SlideTitleColor}{CourseNavy}
\colorlet{SlideSubtitleColor}{CourseGray}
\colorlet{FrameTitleColor}{CourseMaroon}
\colorlet{AccentColor}{CourseBlue}
\colorlet{MutedColor}{CourseGray}
\colorlet{RuleColor}{CourseLine}
\colorlet{FooterTextColor}{CourseGray}

\colorlet{InfoBoxBg}{CourseSky}
\colorlet{InfoBoxDraw}{CourseBlue!28}
\colorlet{SuccessBoxBg}{CourseMint}
\colorlet{SuccessBoxDraw}{CourseGreen!30}
\colorlet{NeutralBoxBg}{CourseLight}
\colorlet{NeutralBoxDraw}{CourseLine}
\colorlet{WarningBoxBg}{CourseOrange!8}
\colorlet{WarningBoxDraw}{CourseOrange!35}

% -------------------- aliases used by cards/boxes --------------------
\colorlet{ThemeBlue}{CourseBlue}
\colorlet{ThemeGreen}{CourseGreen}
\colorlet{ThemeOrange}{CourseOrange}
\colorlet{ThemeNavy}{CourseNavy}
\colorlet{ThemeMaroon}{CourseMaroon}
\colorlet{ThemeBlueLight}{CourseSky}
\colorlet{ThemeGreenLight}{CourseMint}
\colorlet{ThemeOrangeLight}{CourseOrange!8}
\colorlet{ThemeLight}{CourseLight}
\colorlet{ThemeLine}{CourseLine}
\colorlet{ThemeGray}{CourseGray}

\colorlet{ThemeRed}{CourseRed}

% -------------------- typography controls --------------------
\newcommand{\DeckMainFont}{\sffamily}
\newcommand{\DeckTitleSize}{\LARGE}
\newcommand{\DeckSubtitleSize}{\normalsize}
\newcommand{\DeckAuthorSize}{\normalsize}
\newcommand{\DeckInstituteSize}{\small}
\newcommand{\DeckDateSize}{\small}

\newcommand{\FrameTitleSize}{\Large}
\newcommand{\FrameTitleSeries}{\bfseries}
\newcommand{\FrameTitleTopSpace}{0.05em}
\newcommand{\FrameTitleToRuleSpace}{0.25em}
\newcommand{\FrameTitleRuleThickness}{0.45pt}
\newcommand{\FrameTitleAfterRuleSpace}{2pt}

\newcommand{\BodyTextSize}{\normalsize}
\newcommand{\SmallTextSize}{\scriptsize}
\newcommand{\FooterTextSize}{\scriptsize}

% -------------------- global beamer setup --------------------
\setbeamercolor{normal text}{fg=CourseNavy,bg=white}
\setbeamercolor{structure}{fg=AccentColor}

\setbeamercolor{title}{fg=SlideTitleColor}
\setbeamercolor{subtitle}{fg=SlideSubtitleColor}
\setbeamercolor{author}{fg=SlideSubtitleColor}
\setbeamercolor{institute}{fg=CourseNavy}
\setbeamercolor{date}{fg=CourseMaroon}

\setbeamercolor{frametitle}{fg=FrameTitleColor,bg=white}
\setbeamercolor{itemize item}{fg=AccentColor}
\setbeamercolor{itemize subitem}{fg=AccentColor}
\setbeamercolor{enumerate item}{fg=AccentColor}

\setbeamerfont{title}{family=\DeckMainFont,series=\bfseries,size=\DeckTitleSize}
\setbeamerfont{subtitle}{family=\DeckMainFont,size=\DeckSubtitleSize}
\setbeamerfont{author}{family=\DeckMainFont,size=\DeckAuthorSize}
\setbeamerfont{institute}{family=\DeckMainFont,size=\DeckInstituteSize}
\setbeamerfont{date}{family=\DeckMainFont,size=\DeckDateSize}

\setbeamerfont{frametitle}{family=\DeckMainFont,series=\FrameTitleSeries,size=\FrameTitleSize}
\setbeamerfont{normal text}{family=\DeckMainFont,size=\BodyTextSize}
\setbeamerfont{footline}{family=\DeckMainFont,size=\FooterTextSize}

\setbeamertemplate{navigation symbols}{}

% -------------------- frame title style --------------------
\setbeamertemplate{frametitle}{%
  \vskip\FrameTitleTopSpace%
  {\usebeamerfont{frametitle}%
   \usebeamercolor[fg]{frametitle}%
   \insertframetitle\par}%
  \vskip\FrameTitleToRuleSpace%
  \nointerlineskip%
  {\color{RuleColor}%
   \hrule height \FrameTitleRuleThickness width \linewidth}%
  \nointerlineskip%
  \vskip\FrameTitleAfterRuleSpace%
}

% -------------------- footline style --------------------
\newcommand{\FooterLeftText}{}
\setbeamertemplate{footline}{}

% -------------------- reusable text color commands --------------------
\newcommand{\blue}[1]{\textcolor{CourseBlue}{#1}}
\newcommand{\green}[1]{\textcolor{CourseGreen}{#1}}
\newcommand{\orange}[1]{\textcolor{CourseOrange}{#1}}
\newcommand{\muted}[1]{\textcolor{MutedColor}{#1}}

% =========================================================
% CLASSIC BEAMER-STYLE CALLOUT BOXES
% =========================================================

\newcommand{\CalloutTitleFont}{\large\bfseries}
\newcommand{\CalloutBodyFont}{\normalsize}

\newcommand{\CalloutArc}{7pt}
\newcommand{\CalloutRule}{0.55pt}

\newcommand{\CalloutLeftPad}{10pt}
\newcommand{\CalloutRightPad}{10pt}
\newcommand{\CalloutTopPad}{9pt}
\newcommand{\CalloutBottomPad}{9pt}

\newcommand{\CalloutTitleLeftPad}{10pt}
\newcommand{\CalloutTitleRightPad}{10pt}
\newcommand{\CalloutTitleTopPad}{2.5pt}
\newcommand{\CalloutTitleBottomPad}{2.5pt}

\tcbset{
  calloutbase/.style={
    enhanced,
    width=\linewidth,
    colback=white,
    coltitle=white,
    fonttitle=\CalloutTitleFont,
    fontupper=\CalloutBodyFont,
    boxrule=\CalloutRule,
    arc=\CalloutArc,
    left=\CalloutLeftPad,
    right=\CalloutRightPad,
    top=\CalloutTopPad,
    bottom=\CalloutBottomPad,
    drop fuzzy shadow,
    boxed title style={
      sharp corners=south,
      rounded corners=north,
      boxrule=0pt,
      left=\CalloutTitleLeftPad,
      right=\CalloutTitleRightPad,
      top=\CalloutTitleTopPad,
      bottom=\CalloutTitleBottomPad
    },
    attach boxed title to top left={
      xshift=0pt,
      yshift=0pt
    }
  },
  calloutcompacttitle/.style={
    attach boxed title to top left={
      xshift=0pt,
      yshift=0pt
    }
  },
  calloutfulltitle/.style={
    boxed title style={
      width=\linewidth,
      boxrule=0pt,
      sharp corners=south,
      rounded corners=north,
      left=\CalloutTitleLeftPad,
      right=\CalloutTitleRightPad,
      top=\CalloutTitleTopPad,
      bottom=\CalloutTitleBottomPad
    },
    attach boxed title to top left={
      xshift=0pt,
      yshift=0pt
    },
    halign title=left
  },
  calloutwhitebody/.style={colback=white},
  calloutshadedblue/.style={colback=ThemeBlueLight},
  calloutshadedgreen/.style={colback=ThemeGreenLight},
  calloutshadedorange/.style={colback=ThemeOrangeLight},
  calloutshadedgray/.style={colback=ThemeLight}
}

\newtcolorbox{BlueBox}[2][]{
  calloutbase,
  calloutcompacttitle,
  title={#2},
  colframe=ThemeBlue,
  colbacktitle=ThemeBlue,
  #1
}

\newtcolorbox{GreenBox}[2][]{
  calloutbase,
  calloutcompacttitle,
  title={#2},
  colframe=ThemeGreen,
  colbacktitle=ThemeGreen,
  #1
}

\newtcolorbox{OrangeBox}[2][]{
  calloutbase,
  calloutcompacttitle,
  title={#2},
  colframe=ThemeOrange,
  colbacktitle=ThemeOrange,
  #1
}

\newtcolorbox{NavyBox}[2][]{
  calloutbase,
  calloutcompacttitle,
  title={#2},
  colframe=ThemeNavy,
  colbacktitle=ThemeNavy,
  #1
}

\newtcolorbox{PlainBox}[1][]{
  calloutbase,
  title={},
  colframe=ThemeNavy,
  colbacktitle=ThemeNavy,
  #1
}

% =========================================================
% MODERN PRESENTATION CARDS
% =========================================================

\newcommand{\CardArc}{7pt}
\newcommand{\CardRule}{0.45pt}
\newcommand{\CardTopRule}{2.2pt}

\newcommand{\CardLeftPad}{8pt}
\newcommand{\CardRightPad}{8pt}
\newcommand{\CardTopPad}{8pt}
\newcommand{\CardBottomPad}{8pt}

\newcommand{\CardNumberSize}{\Huge}
\newcommand{\CardTitleSize}{\large}
\newcommand{\CardBodySize}{\normalsize}

\tcbset{
  moderncardbase/.style={
    enhanced,
    width=\linewidth,
    colback=white,
    colframe=ThemeLine,
    boxrule=\CardRule,
    arc=\CardArc,
    left=\CardLeftPad,
    right=\CardRightPad,
    top=\CardTopPad,
    bottom=\CardBottomPad,
    drop fuzzy shadow,
    frame hidden=false,
    borderline north={\CardTopRule}{0pt}{#1}
  }
}

\newtcolorbox{BlueCard}[1][]{moderncardbase=ThemeBlue,#1}
\newtcolorbox{GreenCard}[1][]{moderncardbase=ThemeGreen,#1}
\newtcolorbox{OrangeCard}[1][]{moderncardbase=ThemeOrange,#1}
\newtcolorbox{NavyCard}[1][]{moderncardbase=ThemeNavy,#1}

\newcommand{\CardNumber}[2]{%
  {\CardNumberSize\bfseries\textcolor{#1}{#2}}\par
}

\newcommand{\CardIcon}[2]{%
  {\Huge\textcolor{#1}{#2}}\par
}

\newcommand{\CardTitle}[1]{%
  \vspace{0.15cm}
  {\CardTitleSize\bfseries #1}\par
}

\newcommand{\CardText}[1]{%
  \vspace{0.15cm}
  {\CardBodySize\color{ThemeGray} #1}\par
}




% =========================================================
% HORIZONTAL SIDE CARDS
% Left color stripe + horizontal heading + wide content area
% =========================================================

\newenvironment{GreenSideCard}[1]{%
  \begin{tcolorbox}[
    enhanced,
    colback=white,
    colframe=ThemeGreen!65!black,
    boxrule=0.50pt,
    arc=5pt,
    left=10pt,
    right=10pt,
    top=8pt,
    bottom=8pt,
    borderline west={4pt}{0pt}{ThemeGreen},
    drop fuzzy shadow
  ]
  \noindent
  \begin{minipage}[c]{0.2\linewidth}
    \raggedright
    {\Large\bfseries\color{ThemeGreen}#1}
  \end{minipage}%
  \hfill
  \begin{minipage}[c]{0.74\linewidth}
    \raggedright
}{%
  \end{minipage}
  \end{tcolorbox}
}

\newenvironment{OrangeSideCard}[1]{%
  \begin{tcolorbox}[
    enhanced,
    colback=white,
    colframe=ThemeOrange!65!black,
    boxrule=0.50pt,
    arc=5pt,
    left=10pt,
    right=10pt,
    top=8pt,
    bottom=8pt,
    borderline west={4pt}{0pt}{ThemeOrange},
    drop fuzzy shadow
  ]
  \noindent
  \begin{minipage}[c]{0.20\linewidth}
    \raggedright
    {\Large\bfseries\color{ThemeOrange}#1}
  \end{minipage}%
  \hfill
  \begin{minipage}[c]{0.74\linewidth}
    \raggedright
}{%
  \end{minipage}
  \end{tcolorbox}
}

\newenvironment{BlueSideCard}[1]{%
  \begin{tcolorbox}[
    enhanced,
    colback=white,
    colframe=ThemeBlue!65!black,
    boxrule=0.50pt,
    arc=5pt,
    left=10pt,
    right=10pt,
    top=8pt,
    bottom=8pt,
    borderline west={4pt}{0pt}{ThemeBlue},
    drop fuzzy shadow
  ]
  \noindent
  \begin{minipage}[c]{0.20\linewidth}
    \raggedright
    {\Large\bfseries\color{ThemeBlue}#1}
  \end{minipage}%
  \hfill
  \begin{minipage}[c]{0.74\linewidth}
    \raggedright
}{%
  \end{minipage}
  \end{tcolorbox}
}








% =========================================================
% FIGURE PLACEHOLDER HELPERS
% =========================================================

\newcommand{\MissingFigure}[2]{%
\begin{tikzpicture}
  \node[
    rounded corners=8pt,
    draw=CourseLine,
    fill=CourseLight,
    minimum width=#1,
    minimum height=#2,
    align=center,
    inner sep=8pt
  ] {
    \color{CourseGray}
    \faImage\\[4pt]
    Placeholder figure\\
    \scriptsize Replace with your actual file
  };
\end{tikzpicture}%
}

\newcommand{\MaybeFigure}[3]{%
\IfFileExists{#1}{%
  \includegraphics[width=#2,height=#3,keepaspectratio]{#1}%
}{%
  \MissingFigure{#2}{#3}%
}%
}

\newcommand{\SignalPlaceholder}[3]{%
\begin{tcolorbox}[
  enhanced,
  width=\linewidth,
  height=#1,
  colback=#2!8,
  colframe=#2!45,
  boxrule=0.55pt,
  arc=7pt,
  left=5pt,
  right=5pt,
  top=5pt,
  bottom=5pt,
  halign=center,
  valign=center,
  drop fuzzy shadow
]
{\Large\textcolor{#2}{\faImage}}\\[4pt]
{\bfseries #3}\\[-1pt]
{\scriptsize\color{ThemeGray} image placeholder}
\end{tcolorbox}%
}

\newcommand{\TinySignalSketchCT}{%
\begin{tikzpicture}[x=0.62cm,y=0.40cm]
  \draw[-{Latex[length=2mm]},CourseGray] (-2.4,0) -- (2.6,0) node[right] {\scriptsize $t$};
  \draw[-{Latex[length=2mm]},CourseGray] (0,-1.35) -- (0,1.45) node[above] {\scriptsize $x(t)$};
  \draw[very thick,CourseBlue,smooth,domain=-2.1:2.2,samples=80]
    plot (\x,{0.55*sin(150*\x)+0.35*cos(70*\x)});
\end{tikzpicture}%
}

\newcommand{\TinySignalSketchDT}{%
\begin{tikzpicture}[x=0.40cm,y=0.40cm]
  \draw[-{Latex[length=2mm]},CourseGray] (-5.4,0) -- (5.8,0) node[right] {\scriptsize $n$};
  \draw[-{Latex[length=2mm]},CourseGray] (0,-1.4) -- (0,1.5) node[above] {\scriptsize $x[n]$};
  \foreach \n/\v in {-5/-0.4,-4/-0.8,-3/-0.15,-2/0.55,-1/0.9,0/1.2,1/0.75,2/0.35,3/-0.25,4/-0.55,5/-0.15}{
    \draw[very thick,CourseGreen] (\n,0) -- (\n,\v);
    \fill[CourseGreen] (\n,\v) circle (2pt);
  }
\end{tikzpicture}%
}

% =========================================================
% TITLE METADATA
% =========================================================




\title{Lecture 4: Convolution, Causality, and Stability}
\vspace{6cm}
\subtitle{CSE 219: Signals and Linear Systems}

\author{Anik Saha\\[-1pt]{\small Adjunct Lecturer, Department of CSE, BUET}}
% \institute{Department of Computer Science and Engineering\\Bangladesh University of Engineering and Technology}
\date{}

\begin{document}


% =========================================================
% TITLE
% =========================================================

\begin{frame}[plain]
\vspace{0.55cm}
\titlepage

% \vspace{-0.15cm}
% \begin{center}
% \begin{BlueBox}[width=0.78\textwidth,halign=center,calloutshadedblue,fontupper=\large]{}
% A first language for describing information that changes.
% \end{BlueBox}
% \end{center}


\end{frame}


\begin{frame}{Recall: Keep the first signal just as is, flip the second, slide rightward, find overlaps, then multiply-and-add!}

\begin{center}

\begin{figure}[t]
    \centering
    \includegraphics[width=0.85\textwidth]{figs/convolution-penguin-shark.png}
\end{figure}

\end{center}

\end{frame}


% --------------------------------------------------------
\begin{frame}{Recall: discrete-time convolution is sliding multiply-and-add}

\begin{BlueBox}[
  calloutshadedblue,
  fontupper=\large,
  halign=center
]{}
For two discrete-time signals \(x[n]\) and \(h[n]\), their convolution is
\[
\boxed{\displaystyle
y[n]=(x*h)[n]=\sum_{k=-\infty}^{\infty}x[k]\,h[n-k].
}
\]
\end{BlueBox}

\vspace{0.10cm}

\begin{center}
\begin{GreenCard}[
  width=0.92\textwidth,
  halign=center,
  top=10pt,
  bottom=10pt
]
\CardTitle{What the formula is saying}

\vspace{0.06cm}

\begin{tikzpicture}[x=1cm,y=0.90cm,every node/.style={font=\scriptsize}]

\node[
  draw=ThemeBlue,
  fill=ThemeBlue!10,
  rounded corners=5pt,
  very thick,
  minimum width=2.6cm,
  minimum height=0.75cm,
  align=center
] (x) at (-4.2,0) {\Large \(x[k]\)};

\node[
  draw=ThemeOrange,
  fill=ThemeOrange!10,
  rounded corners=5pt,
  very thick,
  minimum width=2.6cm,
  minimum height=0.75cm,
  align=center
] (h) at (0,0) {\Large \(h[n-k]\)};

\node[
  draw=ThemeGreen,
  fill=ThemeGreen!10,
  rounded corners=5pt,
  very thick,
  minimum width=2.7cm,
  minimum height=0.75cm,
  align=center
] (sum) at (4.25,0) {\Large \(y[n]\)};

\draw[-{Latex[length=2.4mm]},very thick,ThemeGray] (x) -- (h)
  node[midway,above] {multiply};
\draw[-{Latex[length=2.4mm]},very thick,ThemeGray] (h) -- (sum)
  node[midway,above] {add over \(k\)};

\node[ThemeBlue] at (0,-1.05)
  {For each fixed \(n\), the sum produces one output sample \(y[n]\).};

\end{tikzpicture}

\end{GreenCard}
\end{center}

\end{frame}






% --------------------------------------------------------
\begin{frame}{Convolution gives a whole output signal, one value at a time}

\begin{center}
\begin{BlueCard}[
  width=\textwidth,
  halign=center,
  top=10pt,
  bottom=10pt
]
% \CardTitle{To compute one output sample, fix \(n\), then sweep over \(k\)}

% \vspace{0.02cm}

\begin{columns}[c,totalwidth=0.96\textwidth]

% ---------------- LEFT PANEL ----------------
\begin{column}{0.5\textwidth}
\begin{GreenCard}[
  halign=center,
  top=8pt,
  bottom=8pt
]
\CardTitle{Sweep over the input index \(k\)}

\vspace{0.04cm}

\begin{tikzpicture}[x=0.44cm,y=0.55cm,every node/.style={font=\scriptsize}]

  % axis
  \draw[-{Latex[length=1.4mm]},thick,black]
    (-6.0,0) -- (6.5,0) node[right] {$k$};

  \draw[-{Latex[length=1.4mm]},thick,black]
    (0,-1.45) -- (0,2.15) node[above] {$x[k]$};

  % continuation dots
  \node[ThemeGray,font=\large] at (-5.45,0.75) {$\cdots$};
  \node[ThemeGray,font=\large] at (5.75,0.75) {$\cdots$};

  % shown samples
  \foreach \k/\v in {-4/0.55,-3/1.10,-2/-0.65,-1/0.90,0/1.55,1/0.45,2/-0.85,3/1.20,4/0.70}{
    \draw[very thick,ThemeGreen] (\k,0) -- (\k,\v);
    \fill[ThemeGreen] (\k,\v) circle (2.0pt);
    \fill[black] (\k,0) circle (0.75pt);
  }

  % ticks
  \foreach \k in {-4,-3,-2,-1,0,1,2,3,4}{
    \draw[black] (\k,0.05) -- (\k,-0.05);
  }

  \node[below] at (-4,-0.06) {$-4$};
  \node[below] at (0,-0.06) {$0$};
  \node[below] at (4,-0.06) {$4$};

  % sweep line
  \draw[{Latex[length=1.6mm]}-{Latex[length=1.6mm]},very thick,ThemeOrange]
    (-4.6,-1.58) -- (4.6,-1.58);

  \node[ThemeOrange,font=\bfseries] at (0,-2.12)
    {sweep over all \(k\)};

\end{tikzpicture}

% \vspace{0.04cm}

% {\normalsize
% For a fixed output location, we inspect all possible input indices \(k\).
% }

\end{GreenCard}
\end{column}

\hfill

% ---------------- RIGHT PANEL ----------------
\begin{column}{0.44\textwidth}
\begin{OrangeCard}[
  halign=center,
  top=8pt,
  bottom=8pt
]
\CardTitle{Fix one output index}

\vspace{0.04cm}

\begin{tikzpicture}[x=0.52cm,y=0.55cm,every node/.style={font=\scriptsize}]

  % output axis
  \draw[-{Latex[length=1.4mm]},thick,black]
    (-3.5,0) -- (6.0,0) node[right] {$n$};

  \draw[-{Latex[length=1.4mm]},thick,black]
    (0,-0.20) -- (0,2.40) node[above] {$y[n]$};

  % gray unknown output samples
  \foreach \n in {-3,-2,-1,0,1,2,4,5}{
    \fill[ThemeGray] (\n,0) circle (1.2pt);
    \draw[black] (\n,0.05) -- (\n,-0.05);
  }

  % highlighted fixed n=3
  \draw[densely dashed,ThemeOrange,very thick] (3,0) -- (3,2.05);

  \draw[very thick,ThemeBlue] (3,0) -- (3,1.45);
  \fill[ThemeBlue] (3,1.45) circle (2.4pt);
  \fill[black] (3,0) circle (0.75pt);

  \node[below] at (0,-0.06) {$0$};
  \node[below] at (3,-0.06) {$3$};
  \node[ThemeOrange,above] at (3,2.08) {fixed \(n=3\)};
  \node[ThemeBlue,right] at (3.18,1.45) {\(y[3]\)};

  \node[ThemeGray,font=\large] at (-3.05,0.75) {$\cdots$};
  \node[ThemeGray,font=\large] at (5.45,0.75) {$\cdots$};

\end{tikzpicture}

% \vspace{0.04cm}

% {\normalsize
% This one highlighted stem is one number: \(y[3]\).
% }

\end{OrangeCard}
\end{column}

\end{columns}

\vspace{0.10cm}

\begin{BlueBox}[
  calloutshadedblue,
  halign=center,
  fontupper=\large
]{}
% For example, when \(n=3\),
\(
y[3]=\sum_{k=-\infty}^{\infty}x[k]\,h[3-k]
\)
\end{BlueBox}

\end{BlueCard}
\end{center}

\end{frame}







% --------------------------------------------------------
\begin{frame}{For fixed \(n\), \(h[n-k]\) is made by reflect-then-right-shift}

% \begin{BlueBox}[
%   calloutshadedblue,
%   fontupper=\large,
%   halign=center
% ]{}
% When computing one output value \(y[n]\), \(n\) is fixed.  
% Then \(h[n-k]\) is viewed as a function of the sweeping variable \(k\).
% \end{BlueBox}

% \vspace{0.08cm}

\begin{center}
\begin{OrangeCard}[
  width=0.94\textwidth,
  halign=center,
  top=9pt,
  bottom=9pt
]
% \CardTitle{From \(h[k]\) to \(h[n-k]\)}

\vspace{0.04cm}

\begin{tikzpicture}[x=0.58cm,y=0.62cm,every node/.style={font=\scriptsize}]

% =========================================================
% Row 1: h[k]
% =========================================================
\begin{scope}[yshift=2.25cm]

  \node[anchor=east,ThemeBlue,font=\large\bfseries] at (-5.15,0.65) {\(h[k]\)};

  \draw[-{Latex[length=1.3mm]},thick,black]
    (-4.6,0) -- (4.9,0) node[right] {$k$};

  \foreach \k in {-4,-3,-2,-1,0,1,2,3,4}{
    \draw[black] (\k,0.05) -- (\k,-0.05);
  }

  \foreach \k/\v in {0/1,1/2,2/1,3/0.5}{
    \draw[very thick,ThemeBlue] (\k,0) -- (\k,\v);
    \fill[ThemeBlue] (\k,\v) circle (2.0pt);
    \fill[black] (\k,0) circle (0.75pt);
  }

  \node[below] at (0,-0.08) {$0$};
  \node[below] at (3,-0.08) {$3$};

\end{scope}


% =========================================================
% Row 2: h[-k]
% =========================================================
\begin{scope}[yshift=0.20cm]

  \node[anchor=east,ThemeOrange,font=\large\bfseries] at (-5.15,0.65) {\(h[-k]\)};

  \draw[-{Latex[length=1.3mm]},thick,black]
    (-4.6,0) -- (4.9,0) node[right] {$k$};

  \foreach \k in {-4,-3,-2,-1,0,1,2,3,4}{
    \draw[black] (\k,0.05) -- (\k,-0.05);
  }

  % reflected version of h[k]
  \foreach \k/\v in {0/1,-1/2,-2/1,-3/0.5}{
    \draw[very thick,ThemeOrange] (\k,0) -- (\k,\v);
    \fill[ThemeOrange] (\k,\v) circle (2.0pt);
    \fill[black] (\k,0) circle (0.75pt);
  }

  \node[below] at (0,-0.08) {$0$};
  \node[below] at (-3,-0.08) {$-3$};

\end{scope}


% =========================================================
% Row 3: h[n-k], example n=2
% =========================================================
\begin{scope}[yshift=-1.85cm]

  \node[anchor=east,ThemeGreen,font=\large\bfseries] at (-5.15,0.65) {\(h[2-k]\)};

  \draw[-{Latex[length=1.3mm]},thick,black]
    (-4.6,0) -- (4.9,0) node[right] {$k$};

  \foreach \k in {-4,-3,-2,-1,0,1,2,3,4}{
    \draw[black] (\k,0.05) -- (\k,-0.05);
  }

  % h[-k] shifted right by 2
  \foreach \k/\v in {2/1,1/2,0/1,-1/0.5}{
    \draw[very thick,ThemeGreen] (\k,0) -- (\k,\v);
    \fill[ThemeGreen] (\k,\v) circle (2.0pt);
    \fill[black] (\k,0) circle (0.75pt);
  }

  \node[below] at (0,-0.08) {$0$};
  \node[below] at (2,-0.08) {$2$};
  \node[ThemeGreen,above] at (5,1.10) {\(h[0]\) now sits at \(k=2\)};

\end{scope}


% =========================================================
% Arrows and annotations
% =========================================================
% \draw[-{Latex[length=2mm]},very thick,ThemeOrange]
%   (5.25,2.90) -- (5.25,1.15);
% \node[ThemeOrange,align=center,font=\bfseries] at (6.35,2.03) {reflect\\around \(0\)};

% \draw[-{Latex[length=2mm]},very thick,ThemeGreen]
%   (5.25,0.85) -- (5.25,-0.95);
% \node[ThemeGreen,align=center,font=\bfseries] at (6.45,-0.05) {shift right\\by \(n=2\)};

\end{tikzpicture}

% \vspace{0.06cm}

% \begin{GreenBox}[
%   calloutshadedgreen,
%   halign=center,
%   fontupper=\large
% ]{}
% In general: \(h[k]\rightarrow h[-k]\rightarrow h[n-k]\).
% \end{GreenBox}

\end{OrangeCard}
\end{center}

\end{frame}



% --------------------------------------------------------
\begin{frame}{The continuous-time version replaces the sum by an integral}

\begin{BlueBox}[
  calloutshadedblue,
  fontupper=\large,
  halign=center
]{}
For continuous-time signals, the idea is the same:  
reflect, shift, multiply, and add continuously.
\end{BlueBox}

\vspace{0.12cm}

\begin{center}
\begin{GreenCard}[
  width=0.90\textwidth,
  halign=center,
  top=12pt,
  bottom=12pt
]
\CardTitle{Continuous-time convolution}

\vspace{0.08cm}

{\Large
\[
\boxed{\displaystyle
y(t)=(x*h)(t)=\int_{-\infty}^{\infty}x(\tau)\,h(t-\tau)\,d\tau
}
\]
}

% \vspace{0.08cm}

% \begin{columns}[T,totalwidth=0.88\textwidth]

% \begin{column}{0.48\textwidth}
% \begin{BlueBox}[
%   calloutshadedblue,
%   halign=center,
%   fontupper=\large
% ]{}
% Discrete time:
% \[
% \sum_k x[k]h[n-k]
% \]
% \end{BlueBox}
% \end{column}

% \begin{column}{0.48\textwidth}
% \begin{OrangeBox}[
%   calloutshadedorange,
%   halign=center,
%   fontupper=\large
% ]{}
% Continuous time:
% \[
% \int x(\tau)h(t-\tau)\,d\tau
% \]
% \end{OrangeBox}
% \end{column}

% \end{columns}

\end{GreenCard}
\end{center}

\end{frame}




% --------------------------------------------------------
\begin{frame}{For a fixed \(t\), sweep along \(\tau\) and integrate along the overlap}

\begin{center}
\begin{BlueCard}[
  width=0.95\textwidth,
  halign=center,
  top=10pt,
  bottom=10pt
]
\CardTitle{Fix \(t\), then slide \(h(t-\tau)\) along the \(\tau\)-axis}

% \vspace{0.06cm}

\begin{tikzpicture}[x=1.00cm,y=1.00cm,every node/.style={font=\scriptsize}]

% ------------------------------------------------
% main axis
% ------------------------------------------------
\draw[-{Latex[length=1.7mm]},thick,black]
  (-3.9,0) -- (5.0,0) node[right] {$\tau$};

\draw[-{Latex[length=1.7mm]},thick,black]
  (0,-0.12) -- (0,1.75);

\foreach \x/\lab in {-3/{},-2/{},-1/{},0/{0},1/{},2/{},3/{},4/{}}{
  \draw[black] (\x,0.04) -- (\x,-0.04);
  \node[below] at (\x,-0.06) {\(\lab\)};
}

% ------------------------------------------------
% x(tau): blue rectangular pulse from -2 to 2
% ------------------------------------------------
\fill[ThemeBlue!12] (-2,0) rectangle (2,1.00);
\draw[very thick,ThemeBlue] (-2,0) -- (-2,1.00) -- (2,1.00) -- (2,0);

\node[ThemeBlue,font=\bfseries] at (-1.15,1.25) {\(x(\tau)\)};

% ------------------------------------------------
% h(t-tau): orange rectangular pulse from 0.8 to 3.8
% ------------------------------------------------
\fill[ThemeOrange!15] (0.8,0) rectangle (3.8,0.72);
\draw[very thick,ThemeOrange] (0.8,0) -- (0.8,0.72) -- (3.8,0.72) -- (3.8,0);

\node[ThemeOrange,font=\bfseries] at (3.00,0.96) {\(h(t-\tau)\)};

% ------------------------------------------------
% overlap region
% ------------------------------------------------
\fill[ThemeGreen!25] (0.8,0) rectangle (2,0.72);
\draw[very thick,ThemeGreen] (0.8,0) rectangle (2,0.72);

\node[ThemeGreen,font=\bfseries] at (1.40,0.37) {overlap};

% small arrows showing sliding
\draw[-{Latex[length=1.9mm]},very thick,ThemeOrange]
  (2.15,-0.45) -- (3.25,-0.45);

\node[ThemeOrange,font=\bfseries] at (2.70,-0.72) {slide as \(t\) changes};

% labels for support endpoints
\node[below,ThemeBlue] at (-2,-0.10) {\(a\)};
\node[below,ThemeBlue] at (2,-0.10) {\(b\)};
\node[below,ThemeOrange] at (0.8,-0.10) {\(t-d\)};
\node[below,ThemeOrange] at (3.8,-0.10) {\(t-c\)};

\end{tikzpicture}

\vspace{0.10cm}

\begin{GreenBox}[
  calloutshadedgreen,
  halign=center,
  fontupper=\large
]{}
For this one fixed \(t\),
\[
y(t)=\int_{-\infty}^{\infty} x(\tau)\,h(t-\tau)\,d\tau
\]
% is the area under the product \(x(\tau)h(t-\tau)\).
\end{GreenBox}

\end{BlueCard}
\end{center}

\end{frame}



% --------------------------------------------------------
\begin{frame}{Discrete and continuous convolution follow the same recipe}

\begin{center}

\begin{GreenSideCard}{1}
{\large\bfseries Choose the output location.}

\vspace{0.03cm}

{\normalsize
For discrete time, choose \(n\).  
For continuous time, choose \(t\).
}
\end{GreenSideCard}

\vspace{0.10cm}

\begin{OrangeSideCard}{2}
{\large\bfseries Reflect and shift the second signal.}

\vspace{0.03cm}

{\normalsize
Discrete: \(h[k]\to h[-k]\to h[n-k]\).  
Continuous: \(h(\tau)\to h(-\tau)\to h(t-\tau)\).
}
\end{OrangeSideCard}

\vspace{0.10cm}

\begin{BlueSideCard}{3}
{\large\bfseries Multiply where they overlap, then add.}

\vspace{0.03cm}

{\normalsize
Discrete: sum over \(k\).  
Continuous: integrate over \(\tau\).
}
\end{BlueSideCard}

% \vspace{0.12cm}

% \begin{BlueBox}[
%   calloutshadedblue,
%   halign=center,
%   fontupper=\large
% ]{}
% Repeating this for all \(n\) or all \(t\) gives the whole output signal.
% \end{BlueBox}

\end{center}

\end{frame}





% --------------------------------------------------------
\begin{frame}[plain]

% \vfill

\begin{center}

\begin{BlueCard}[
  width=\textwidth,
  halign=center,
  top=12pt,
  bottom=12pt
]

{\Large\bfseries Please try this beautiful interactive convolution demo!}\par

\vspace{0.24cm}

{\large
\href{https://phiresky.github.io/convolution-demo/}
{\texttt{https://phiresky.github.io/convolution-demo/}}
\par
}

\vspace{0.24cm}

\href{https://phiresky.github.io/convolution-demo/}{%
  \includegraphics[width=0.75\textwidth]{figs/convolution-demo.png}
}

% \vspace{0.14cm}

% \begin{OrangeBox}[
%   calloutshadedorange,
%   halign=center,
%   fontupper=\large
% ]{}
% Click the link or the screenshot to open the demo.
% \end{OrangeBox}

\end{BlueCard}

\end{center}

\vfill

\end{frame}





% --------------------------------------------------------
\begin{frame}[plain]

\vfill

\begin{center}

{\fontsize{30}{34}\selectfont\bfseries\color{FrameTitleColor}
Practice Problems on Convolution\par}

% \vspace{0.25cm}

% {\Large\color{ThemeGray}
% Convolution is simply a mathematical operation that appears very universally.\par}

\end{center}

\vfill

\end{frame}





% --------------------------------------------------------
\begin{frame}{Practice (Example 2.2)}

\begin{OrangeBox}[
  calloutshadedorange,
  fontupper=\large,
  halign=center
]{}
For the two signals below, compute
\[
y[n]=x[n]*h[n]=\sum_{k=-\infty}^{\infty}x[k]h[n-k].
\]
\end{OrangeBox}

\vspace{0.02cm}

\begin{columns}[c,totalwidth=\textwidth]

% x[k]
\begin{column}{0.48\textwidth}
\begin{GreenCard}[halign=center,top=10pt,bottom=10pt]
\CardTitle{Input \(x[k]\)}

\vspace{0.04cm}

\begin{tikzpicture}[x=0.72cm,y=0.55cm,every node/.style={font=\scriptsize}]
  \draw[-{Latex[length=1.5mm]},thick,black] (-3.2,0) -- (4.2,0) node[right] {$k$};
  \draw[-{Latex[length=1.5mm]},thick,black] (0,-0.15) -- (0,2.45);

  \foreach \k in {-3,-2,-1,0,1,2,3}{
    \fill[black] (\k,0) circle (0.8pt);
    \draw[black] (\k,0.04) -- (\k,-0.04);
    \node[below] at (\k,-0.06) {$\k$};
  }

  \draw[very thick,ThemeGreen] (0,0) -- (0,0.5);
  \fill[ThemeGreen] (0,0.5) circle (2.1pt);
  \node[left] at (0,0.5) {\(0.5\)};

  \draw[very thick,ThemeGreen] (1,0) -- (1,2);
  \fill[ThemeGreen] (1,2) circle (2.1pt);
  \node[above] at (1,2.05) {\(2\)};

  \node[ThemeGreen,font=\bfseries] at (2.7,1.60) {\(x[k]\)};
\end{tikzpicture}

% \[
% x[0]=0.5,\qquad x[1]=2.
% \]

\end{GreenCard}
\end{column}

% h[k]
\begin{column}{0.48\textwidth}
\begin{BlueCard}[halign=center,top=10pt,bottom=10pt]
\CardTitle{Impulse response \(h[k]\)}

\vspace{0.04cm}

\begin{tikzpicture}[x=0.72cm,y=0.55cm,every node/.style={font=\scriptsize}]
  \draw[-{Latex[length=1.5mm]},thick,black] (-3.2,0) -- (4.2,0) node[right] {$k$};
  \draw[-{Latex[length=1.5mm]},thick,black] (0,-0.15) -- (0,1.55);

  \foreach \k in {-3,-2,-1,0,1,2,3}{
    \fill[black] (\k,0) circle (0.8pt);
    \draw[black] (\k,0.04) -- (\k,-0.04);
    \node[below] at (\k,-0.06) {$\k$};
  }

  \foreach \k in {0,1,2}{
    \draw[very thick,ThemeBlue] (\k,0) -- (\k,1);
    \fill[ThemeBlue] (\k,1) circle (2.1pt);
  }

  \node[left] at (0,1) {\(1\)};
  \node[ThemeBlue,font=\bfseries] at (2.9,1.18) {\(h[k]\)};
\end{tikzpicture}

% \[
% h[0]=h[1]=h[2]=1.
% \]

\end{BlueCard}
\end{column}

\end{columns}

\vspace{0.10cm}

\begin{BlueBox}[calloutshadedblue,halign=center,fontupper=\large]{}
Use the sliding view: reflect \(h[k]\), shift it to \(h[n-k]\), multiply with \(x[k]\), then add.
\end{BlueBox}

\end{frame}







% --------------------------------------------------------
\begin{frame}{Answer setup: \(h[n-k]\) is the reflected signal shifted by \(n\)}

\begin{BlueBox}[
  calloutshadedblue,
  fontupper=\large,
  halign=center
]{}
For each fixed \(n\), draw \(h[n-k]\) as a function of \(k\).
\end{BlueBox}

\vspace{0.10cm}

\begin{center}
\begin{OrangeCard}[width=0.94\textwidth,halign=center,top=10pt,bottom=10pt]
% \CardTitle{Reflect first, then slide}

% \vspace{0.06cm}

\begin{tikzpicture}[x=0.72cm,y=0.55cm,every node/.style={font=\scriptsize}]

% ---------------- h[k]
\begin{scope}[yshift=2.35cm]
  \node[anchor=east,ThemeBlue,font=\bfseries] at (-4.0,0.55) {\(h[k]\)};
  \draw[-{Latex[length=1.3mm]},thick,black] (-3.4,0) -- (4.0,0) node[right] {$k$};

  \foreach \k in {-2,-1,0,1,2,3}{
    \fill[black] (\k,0) circle (0.75pt);
    \node[below] at (\k,-0.06) {$\k$};
  }

  \foreach \k in {0,1,2}{
    \draw[very thick,ThemeBlue] (\k,0) -- (\k,1);
    \fill[ThemeBlue] (\k,1) circle (2pt);
  }
\end{scope}

% ---------------- h[-k]
\begin{scope}[yshift=0.45cm]
  \node[anchor=east,ThemeOrange,font=\bfseries] at (-4.0,0.55) {\(h[-k]\)};
  \draw[-{Latex[length=1.3mm]},thick,black] (-3.4,0) -- (4.0,0) node[right] {$k$};

  \foreach \k in {-2,-1,0,1,2,3}{
    \fill[black] (\k,0) circle (0.75pt);
    \node[below] at (\k,-0.06) {$\k$};
  }

  \foreach \k in {-2,-1,0}{
    \draw[very thick,ThemeOrange] (\k,0) -- (\k,1);
    \fill[ThemeOrange] (\k,1) circle (2pt);
  }

  \draw[-{Latex[length=1.6mm]},very thick,ThemeOrange] (1.4,0.65) -- (2.35,0.65);
  \node[ThemeOrange,above] at (1.88,0.74) {shift by \(n\)};
\end{scope}

% ---------------- h[n-k]
\begin{scope}[yshift=-1.45cm]
  \node[anchor=east,ThemeGreen,font=\bfseries] at (-4.0,0.55) {\(h[n-k]\)};
  \draw[-{Latex[length=1.3mm]},thick,black] (-3.4,0) -- (4.0,0) node[right] {$k$};

  \foreach \k in {-2,-1,0,1,2,3}{
    \fill[black] (\k,0) circle (0.75pt);
  }

  \draw[very thick,ThemeGreen] (0.55,0) -- (0.55,1);
  \draw[very thick,ThemeGreen] (1.55,0) -- (1.55,1);
  \draw[very thick,ThemeGreen] (2.55,0) -- (2.55,1);
  \fill[ThemeGreen] (0.55,1) circle (2pt);
  \fill[ThemeGreen] (1.55,1) circle (2pt);
  \fill[ThemeGreen] (2.55,1) circle (2pt);

  \node[below] at (0.55,-0.06) {\(n-2\)};
  \node[below] at (1.55,-0.06) {\(n-1\)};
  \node[below] at (2.55,-0.06) {\(n\)};
\end{scope}

\end{tikzpicture}

\vspace{0.04cm}

\[
h[n-k]\neq 0
\quad\text{only at}\quad
k=n-2,\;n-1,\;n.
\]

\end{OrangeCard}
\end{center}

\end{frame}




% --------------------------------------------------------
\begin{frame}{Answer: start sliding \(h[n-k]\) past \(x[k]\)}

\begin{center}

% ---------- n < 0 ----------
\begin{BlueCard}[width=0.94\textwidth,halign=center,top=7pt,bottom=7pt]
\CardTitle{\(n<0\): no overlap}

\vspace{0.02cm}

\begin{tikzpicture}[x=0.58cm,y=0.52cm,every node/.style={font=\scriptsize}]

% ---------------- top axis: x[k]
\draw[-{Latex[length=1.2mm]},thick,black] (-4.7,1.35) -- (4.7,1.35) node[right] {$k$};
\node[anchor=east,ThemeGreen,font=\bfseries] at (-4.9,1.85) {\(x[k]\)};

\foreach \k in {-4,-3,-2,-1,0,1,2,3,4}{
  \fill[black] (\k,1.35) circle (0.65pt);
}

\draw[very thick,ThemeGreen] (0,1.35) -- (0,1.80);
\fill[ThemeGreen] (0,1.80) circle (1.8pt);
\node[above] at (0,1.85) {\(0.5\)};

\draw[very thick,ThemeGreen] (1,1.35) -- (1,2.55);
\fill[ThemeGreen] (1,2.55) circle (1.8pt);
\node[above] at (1,2.60) {\(2\)};

% ---------------- bottom axis: h[n-k]
\draw[-{Latex[length=1.2mm]},thick,black] (-4.7,-0.75) -- (4.7,-0.75) node[right] {$k$};
\node[anchor=east,ThemeOrange,font=\bfseries] at (-4.9,-0.25) {\(h[n-k]\)};

\foreach \k in {-4,-3,-2,-1,0,1,2,3,4}{
  \fill[black] (\k,-0.75) circle (0.65pt);
  \node[below] at (\k,-0.84) {$\k$};
}

% example n=-1, so positions are -3,-2,-1
\foreach \k in {-3,-2,-1}{
  \draw[very thick,ThemeOrange] (\k,-0.75) -- (\k,0.05);
  \fill[ThemeOrange] (\k,0.05) circle (1.8pt);
}

\node[ThemeBlue,font=\large\bfseries] at (5.80,0.55) {\(y[n]=0\)};
% \node[ThemeGray] at (0,-2.35) {for \(n<0\), the two signals do not share any nonzero overlap of \(k\)};

\end{tikzpicture}

For \(n<0\), the two signals do not share any nonzero overlap of \(k\)

\end{BlueCard}

\end{center}

\end{frame}







% --------------------------------------------------------
\begin{frame}{Answer: start sliding \(h[n-k]\) past \(x[k]\)}

\begin{center}

% ---------- n = 0 ----------
\begin{GreenCard}[width=0.94\textwidth,halign=center,top=7pt,bottom=7pt]
\CardTitle{\(n=0\): one overlap}

\vspace{0.02cm}

\begin{tikzpicture}[x=0.58cm,y=0.52cm,every node/.style={font=\scriptsize}]

% ---------------- top axis: x[k]
\draw[-{Latex[length=1.2mm]},thick,black] (-4.7,1.35) -- (4.7,1.35) node[right] {$k$};
\node[anchor=east,ThemeGreen,font=\bfseries] at (-4.9,1.85) {\(x[k]\)};

\foreach \k in {-4,-3,-2,-1,0,1,2,3,4}{
  \fill[black] (\k,1.35) circle (0.65pt);
}

\draw[very thick,ThemeGreen] (0,1.35) -- (0,1.80);
\fill[ThemeGreen] (0,1.80) circle (1.8pt);
\node[above] at (0,1.85) {\(0.5\)};

\draw[very thick,ThemeGreen] (1,1.35) -- (1,2.55);
\fill[ThemeGreen] (1,2.55) circle (1.8pt);
\node[above] at (1,2.60) {\(2\)};

% ---------------- bottom axis: h[0-k]
\draw[-{Latex[length=1.2mm]},thick,black] (-4.7,-0.75) -- (4.7,-0.75) node[right] {$k$};
\node[anchor=east,ThemeOrange,font=\bfseries] at (-4.9,-0.25) {\(h[0-k]\)};

\foreach \k in {-4,-3,-2,-1,0,1,2,3,4}{
  \fill[black] (\k,-0.75) circle (0.65pt);
  \node[below] at (\k,-0.84) {$\k$};
}

% n=0, positions -2,-1,0
\foreach \k in {-2,-1,0}{
  \draw[very thick,ThemeOrange] (\k,-0.75) -- (\k,0.05);
  \fill[ThemeOrange] (\k,0.05) circle (1.8pt);
}

% overlap guide
\draw[densely dashed,very thick,ThemeBlue] (0,1.80) -- (0,0.05);
\node[ThemeBlue,font=\bfseries] at (0.55,0.55) {overlap};

\node[ThemeBlue,font=\large\bfseries] at (5.85,0.55) {\(y[0]=0.5 \times1 = 0.5\)};

\end{tikzpicture}
\end{GreenCard}

\end{center}

\end{frame}

















% --------------------------------------------------------
\begin{frame}{Answer: when the signals overlap more, the sum grows}

\begin{center}

% ---------- n = 1 ----------
\begin{OrangeCard}[width=0.94\textwidth,halign=center,top=7pt,bottom=7pt]
\CardTitle{\(n=1\): two overlaps}

\vspace{0.02cm}

\begin{tikzpicture}[x=0.58cm,y=0.52cm,every node/.style={font=\scriptsize}]

% ---------------- top axis: x[k]
\draw[-{Latex[length=1.2mm]},thick,black] (-4.7,1.35) -- (4.7,1.35) node[right] {$k$};
\node[anchor=east,ThemeGreen,font=\bfseries] at (-4.9,1.85) {\(x[k]\)};

\foreach \k in {-4,-3,-2,-1,0,1,2,3,4}{
  \fill[black] (\k,1.35) circle (0.65pt);
}

\draw[very thick,ThemeGreen] (0,1.35) -- (0,1.80);
\fill[ThemeGreen] (0,1.80) circle (1.8pt);
\node[above] at (0,1.85) {\(0.5\)};

\draw[very thick,ThemeGreen] (1,1.35) -- (1,2.55);
\fill[ThemeGreen] (1,2.55) circle (1.8pt);
\node[above] at (1,2.60) {\(2\)};

% ---------------- bottom axis: h[1-k]
\draw[-{Latex[length=1.2mm]},thick,black] (-4.7,-0.75) -- (4.7,-0.75) node[right] {$k$};
\node[anchor=east,ThemeOrange,font=\bfseries] at (-4.9,-0.25) {\(h[1-k]\)};

\foreach \k in {-4,-3,-2,-1,0,1,2,3,4}{
  \fill[black] (\k,-0.75) circle (0.65pt);
  \node[below] at (\k,-0.84) {$\k$};
}

% n=1, positions -1,0,1
\foreach \k in {-1,0,1}{
  \draw[very thick,ThemeOrange] (\k,-0.75) -- (\k,0.05);
  \fill[ThemeOrange] (\k,0.05) circle (1.8pt);
}

% overlap guides
\foreach \k in {0,1}{
  \draw[densely dashed,very thick,ThemeBlue] (\k,1.80) -- (\k,0.05);
}

\node[ThemeBlue,font=\large\bfseries] at (6.05,0.55) {\(y[1]=0.5+2=2.5\)};

\end{tikzpicture}
\end{OrangeCard}

\end{center}

\end{frame}







% --------------------------------------------------------
\begin{frame}{Answer: when the signals overlap more, the sum grows}

\begin{center}


% ---------- n = 2 ----------
\begin{OrangeCard}[width=0.94\textwidth,halign=center,top=7pt,bottom=7pt]
\CardTitle{\(n=2\): still two overlaps}

\vspace{0.02cm}

\begin{tikzpicture}[x=0.58cm,y=0.52cm,every node/.style={font=\scriptsize}]

% ---------------- top axis: x[k]
\draw[-{Latex[length=1.2mm]},thick,black] (-4.7,1.35) -- (4.7,1.35) node[right] {$k$};
\node[anchor=east,ThemeGreen,font=\bfseries] at (-4.9,1.85) {\(x[k]\)};

\foreach \k in {-4,-3,-2,-1,0,1,2,3,4}{
  \fill[black] (\k,1.35) circle (0.65pt);
}

\draw[very thick,ThemeGreen] (0,1.35) -- (0,1.80);
\fill[ThemeGreen] (0,1.80) circle (1.8pt);
\node[above] at (0,1.85) {\(0.5\)};

\draw[very thick,ThemeGreen] (1,1.35) -- (1,2.55);
\fill[ThemeGreen] (1,2.55) circle (1.8pt);
\node[above] at (1,2.60) {\(2\)};

% ---------------- bottom axis: h[2-k]
\draw[-{Latex[length=1.2mm]},thick,black] (-4.7,-0.75) -- (4.7,-0.75) node[right] {$k$};
\node[anchor=east,ThemeOrange,font=\bfseries] at (-4.9,-0.25) {\(h[2-k]\)};

\foreach \k in {-4,-3,-2,-1,0,1,2,3,4}{
  \fill[black] (\k,-0.75) circle (0.65pt);
  \node[below] at (\k,-0.84) {$\k$};
}

% n=2, positions 0,1,2
\foreach \k in {0,1,2}{
  \draw[very thick,ThemeOrange] (\k,-0.75) -- (\k,0.05);
  \fill[ThemeOrange] (\k,0.05) circle (1.8pt);
}

% overlap guides
\foreach \k in {0,1}{
  \draw[densely dashed,very thick,ThemeBlue] (\k,1.80) -- (\k,0.05);
}

\node[ThemeBlue,font=\large\bfseries] at (6.05,0.55) {\(y[2]=0.5+2=2.5\)};

\end{tikzpicture}
\end{OrangeCard}

\end{center}

\end{frame}













% --------------------------------------------------------
\begin{frame}{Answer: after the last overlap, the output returns to zero}

\begin{center}

% ---------- n = 3 ----------
\begin{GreenCard}[width=0.94\textwidth,halign=center,top=7pt,bottom=7pt]
\CardTitle{\(n=3\): one overlap}

\vspace{0.02cm}

\begin{tikzpicture}[x=0.58cm,y=0.52cm,every node/.style={font=\scriptsize}]

% ---------------- top axis: x[k]
\draw[-{Latex[length=1.2mm]},thick,black] (-4.7,1.35) -- (4.7,1.35) node[right] {$k$};
\node[anchor=east,ThemeGreen,font=\bfseries] at (-4.9,1.85) {\(x[k]\)};

\foreach \k in {-4,-3,-2,-1,0,1,2,3,4}{
  \fill[black] (\k,1.35) circle (0.65pt);
}

\draw[very thick,ThemeGreen] (0,1.35) -- (0,1.80);
\fill[ThemeGreen] (0,1.80) circle (1.8pt);
\node[above] at (0,1.85) {\(0.5\)};

\draw[very thick,ThemeGreen] (1,1.35) -- (1,2.55);
\fill[ThemeGreen] (1,2.55) circle (1.8pt);
\node[above] at (1,2.60) {\(2\)};

% ---------------- bottom axis: h[3-k]
\draw[-{Latex[length=1.2mm]},thick,black] (-4.7,-0.75) -- (4.7,-0.75) node[right] {$k$};
\node[anchor=east,ThemeOrange,font=\bfseries] at (-4.9,-0.25) {\(h[3-k]\)};

\foreach \k in {-4,-3,-2,-1,0,1,2,3,4}{
  \fill[black] (\k,-0.75) circle (0.65pt);
  \node[below] at (\k,-0.84) {$\k$};
}

% n=3, positions 1,2,3
\foreach \k in {1,2,3}{
  \draw[very thick,ThemeOrange] (\k,-0.75) -- (\k,0.05);
  \fill[ThemeOrange] (\k,0.05) circle (1.8pt);
}

% overlap guide
\draw[densely dashed,very thick,ThemeBlue] (1,2.55) -- (1,0.05);

\node[ThemeBlue,font=\large\bfseries] at (5.85,0.55) {\(y[3]=2\)};

\end{tikzpicture}
\end{GreenCard}

\end{center}

\end{frame}












% --------------------------------------------------------
\begin{frame}{Answer: after the last overlap, the output returns to zero}

\begin{center}

% ---------- n > 3 ----------
\begin{BlueCard}[width=0.94\textwidth,halign=center,top=7pt,bottom=7pt]
\CardTitle{\(n>3\): no overlap again}

\vspace{0.02cm}

\begin{tikzpicture}[x=0.58cm,y=0.52cm,every node/.style={font=\scriptsize}]

% ---------------- top axis: x[k]
\draw[-{Latex[length=1.2mm]},thick,black] (-4.7,1.35) -- (5.4,1.35) node[right] {$k$};
\node[anchor=east,ThemeGreen,font=\bfseries] at (-4.9,1.85) {\(x[k]\)};

\foreach \k in {-4,-3,-2,-1,0,1,2,3,4,5}{
  \fill[black] (\k,1.35) circle (0.65pt);
}

\draw[very thick,ThemeGreen] (0,1.35) -- (0,1.80);
\fill[ThemeGreen] (0,1.80) circle (1.8pt);
\node[above] at (0,1.85) {\(0.5\)};

\draw[very thick,ThemeGreen] (1,1.35) -- (1,2.55);
\fill[ThemeGreen] (1,2.55) circle (1.8pt);
\node[above] at (1,2.60) {\(2\)};

% ---------------- bottom axis: h[n-k]
\draw[-{Latex[length=1.2mm]},thick,black] (-4.7,-0.75) -- (5.4,-0.75) node[right] {$k$};
\node[anchor=east,ThemeOrange,font=\bfseries] at (-4.9,-0.25) {\(h[n-k]\)};

\foreach \k in {-4,-3,-2,-1,0,1,2,3,4,5}{
  \fill[black] (\k,-0.75) circle (0.65pt);
  \node[below] at (\k,-0.84) {$\k$};
}

% example n=4, positions 2,3,4
\foreach \k in {2,3,4}{
  \draw[very thick,ThemeOrange] (\k,-0.75) -- (\k,0.05);
  \fill[ThemeOrange] (\k,0.05) circle (1.8pt);
}

\node[ThemeBlue,font=\large\bfseries] at (6.25,0.55) {\(y[n]=0\)};
% \node[ThemeGray] at (0,-1.55) {after \(n=3\), the sliding signal has passed the support of \(x[k]\)};

\end{tikzpicture}

After \(n=3\), the sliding signal has passed the support of \(x[k]\)

\end{BlueCard}

\end{center}

\end{frame}




















% --------------------------------------------------------
\begin{frame}{Answer: final convolution output}

\begin{center}
\begin{BlueCard}[width=0.92\textwidth,halign=center,top=10pt,bottom=10pt]
% \CardTitle{Final convolution output}

\vspace{0.06cm}

\(
y[n]=
\begin{cases}
0.5, & n=0,\\
2.5, & n=1,\\
2.5, & n=2,\\
2, & n=3,\\
0, & \text{otherwise.}
\end{cases}
\)

\vspace{0.04cm}

\begin{tikzpicture}[x=0.78cm,y=0.55cm,every node/.style={font=\scriptsize}]
  \draw[-{Latex[length=1.5mm]},thick,black] (-2.6,0) -- (5.2,0) node[right] {$n$};
  \draw[-{Latex[length=1.5mm]},thick,black] (0,-0.15) -- (0,3.05) node[above] {$y[n]$};

  \foreach \n in {-2,-1,0,1,2,3,4}{
    \fill[black] (\n,0) circle (0.8pt);
    \draw[black] (\n,0.04) -- (\n,-0.04);
    \node[below] at (\n,-0.06) {$\n$};
  }

  \foreach \n/\v/\lab in {0/0.5/{0.5},1/2.5/{2.5},2/2.5/{2.5},3/2/{2}}{
    \draw[very thick,ThemeBlue] (\n,0) -- (\n,\v);
    \fill[ThemeBlue] (\n,\v) circle (2.2pt);
    \node[above] at (\n,\v+0.05) {\(\lab\)};
  }

  \node[left] at (0,1) {$1$};
  \node[left] at (0,2) {$2$};
\end{tikzpicture}

% \vspace{0.06cm}

% \begin{OrangeBox}[calloutshadedorange,halign=center,fontupper=\large]{}
% Each value of \(y[n]\) came from one overlap position of the sliding signal \(h[n-k]\).
% \end{OrangeBox}

\end{BlueCard}
\end{center}

\end{frame}






% =========================================================
% Controls for the exponential-step convolution example
% Change these two values if needed.
% =========================================================
\newcommand{\ExpConvAlpha}{0.82}   % plotted value of alpha
\newcommand{\ExpConvMax}{10}       % number of right-side samples shown
\newcommand{\ExpConvN}{8}          % example fixed n used in overlap pictures



% --------------------------------------------------------
\begin{frame}{Practice (Example 2.3)}

\begin{OrangeBox}[
  calloutshadedorange,
  fontupper=\large,
  halign=center
]{}
Let \(0<\alpha<1\). Compute
\[
y[n]=x[n]*h[n],
\qquad
x[n]=\alpha^n u[n],
\qquad
h[n]=u[n].
\]
\end{OrangeBox}

\vspace{0.10cm}

\begin{columns}[c,totalwidth=\textwidth]

% x[n]
\begin{column}{0.49\textwidth}
\begin{GreenCard}[halign=center,top=10pt,bottom=10pt]
\CardTitle{Input \(x[n]=\alpha^n u[n]\)}

\vspace{0.04cm}

\begin{tikzpicture}[x=0.35cm,y=1.25cm,every node/.style={font=\scriptsize}]

  \draw[-{Latex[length=1.4mm]},thick,black]
    (-5.2,0) -- (16.2,0) node[right] {$n$};
  \draw[-{Latex[length=1.4mm]},thick,black]
    (0,-0.10) -- (0,1.25);

  \node[ThemeGray,font=\large] at (-4.55,0.20) {$\cdots$};
  \node[ThemeGray,font=\large] at (15.45,0.20) {$\cdots$};

  % zero samples before 0
  \foreach \n in {-4,-3,-2,-1}{
    \fill[black] (\n,0) circle (1.0pt);
  }

  % alpha^n samples
  \foreach \n in {0,...,\ExpConvMax}{
    \pgfmathsetmacro{\v}{pow(\ExpConvAlpha,\n)}
    \draw[very thick,ThemeGreen] (\n,0) -- (\n,\v);
    \fill[ThemeGreen] (\n,\v) circle (1.7pt);
    \fill[black] (\n,0) circle (0.65pt);
  }

  \foreach \n in {-4,-3,-2,-1,0,1,2,3,4,5,6,7,8,9,10,11,12,13,14}{
    \draw[black] (\n,0.025) -- (\n,-0.025);
  }

  \node[below] at (0,-0.04) {$0$};
  \node[below] at (5,-0.04) {$5$};
  \node[below] at (10,-0.04) {$10$};
  \node[left] at (0,1) {$1$};

  \node[ThemeGreen,above] at (6.5,0.48) {\(x[n]\)};
\end{tikzpicture}

% \vspace{0.04cm}

% {\normalsize plotted with \(\alpha=\ExpConvAlpha\)}

\end{GreenCard}
\end{column}

% h[n]
\begin{column}{0.49\textwidth}
\begin{BlueCard}[halign=center,top=10pt,bottom=10pt]
\CardTitle{Impulse response \(h[n]=u[n]\)}

\vspace{0.04cm}

\begin{tikzpicture}[x=0.35cm,y=1.25cm,every node/.style={font=\scriptsize}]

  \draw[-{Latex[length=1.4mm]},thick,black]
    (-5.2,0) -- (16.2,0) node[right] {$n$};
  \draw[-{Latex[length=1.4mm]},thick,black]
    (0,-0.10) -- (0,1.25);

  \node[ThemeGray,font=\large] at (-4.55,0.20) {$\cdots$};
  \node[ThemeGray,font=\large] at (15.45,1.08) {$\cdots$};

  % zero samples before 0
  \foreach \n in {-4,-3,-2,-1}{
    \fill[black] (\n,0) circle (1.0pt);
  }

  % unit step samples
  \foreach \n in {0,...,\ExpConvMax}{
    \draw[very thick,ThemeBlue] (\n,0) -- (\n,1);
    \fill[ThemeBlue] (\n,1) circle (1.7pt);
    \fill[black] (\n,0) circle (0.65pt);
  }

  \foreach \n in {-4,-3,-2,-1,0,1,2,3,4,5,6,7,8,9,10,11,12,13,14}{
    \draw[black] (\n,0.025) -- (\n,-0.025);
  }

  \node[below] at (0,-0.04) {$0$};
  \node[below] at (5,-0.04) {$5$};
  \node[below] at (10,-0.04) {$10$};
  \node[left] at (0,1) {$1$};

  \node[ThemeBlue,above] at (8.5,1.04) {\(h[n]\)};

\end{tikzpicture}

% \vspace{0.04cm}

% {\normalsize a unit step impulse response}

\end{BlueCard}
\end{column}

\end{columns}

\end{frame}







% --------------------------------------------------------
\begin{frame}{First understand how \(h[n-k]\) looks like}

% \begin{BlueBox}[
%   calloutshadedblue,
%   fontupper=\large,
%   halign=center
% ]{}
% Since \(h[n]=u[n]\), we have \(h[n-k]=u[n-k]\).  
% So \(h[n-k]=1\) exactly when \(k\le n\).
% \end{BlueBox}

% \vspace{0.10cm}

\begin{center}
\begin{OrangeCard}[width=0.94\textwidth,halign=center,top=10pt,bottom=10pt]
\CardTitle{Reflect \(h[k]\), then shift it by \(n\)}

\vspace{0.05cm}

\begin{tikzpicture}[x=0.42cm,y=0.65cm,every node/.style={font=\scriptsize}]

% ---------------- h[k]
\begin{scope}[yshift=2.25cm]
  \node[anchor=east,ThemeBlue,font=\large\bfseries] at (-6.0,0.55) {\(h[k]\)};
  \draw[-{Latex[length=1.3mm]},thick,black] (-5.5,0) -- (13.3,0) node[right] {$k$};

  \node[ThemeGray,font=\large] at (-4.9,0.25) {$\cdots$};
  \node[ThemeGray,font=\large] at (12.55,1.05) {$\cdots$};

  \foreach \k in {-4,-3,-2,-1}{
    \fill[black] (\k,0) circle (0.75pt);
  }

  \foreach \k in {0,...,12}{
    \draw[very thick,ThemeBlue] (\k,0) -- (\k,1);
    \fill[ThemeBlue] (\k,1) circle (1.6pt);
    \fill[black] (\k,0) circle (0.6pt);
  }

  \node[below] at (0,-0.05) {$0$};
\end{scope}

% ---------------- h[-k]
\begin{scope}[yshift=0.25cm]
  \node[anchor=east,ThemeOrange,font=\large\bfseries] at (-6.0,0.55) {\(h[-k]\)};
  \draw[-{Latex[length=1.3mm]},thick,black] (-12.5,0) -- (6.0,0) node[right] {$k$};

  \node[ThemeGray,font=\large] at (-11.8,1.05) {$\cdots$};
  \node[ThemeGray,font=\large] at (5.25,0.25) {$\cdots$};

  \foreach \k in {-12,...,0}{
    \draw[very thick,ThemeOrange] (\k,0) -- (\k,1);
    \fill[ThemeOrange] (\k,1) circle (1.6pt);
    \fill[black] (\k,0) circle (0.6pt);
  }

  \foreach \k in {1,2,3,4,5}{
    \fill[black] (\k,0) circle (0.75pt);
  }

  \node[below] at (0,-0.05) {$0$};
\end{scope}

% ---------------- h[n-k]
\begin{scope}[yshift=-1.75cm]
  \node[anchor=east,ThemeGreen,font=\large\bfseries] at (-6.0,0.55) {\(h[n-k]\)};
  \draw[-{Latex[length=1.3mm]},thick,black] (-7.0,0) -- (13.3,0) node[right] {$k$};

  \node[ThemeGray,font=\large] at (-6.35,1.05) {$\cdots$};
  \node[ThemeGray,font=\large] at (12.55,0.25) {$\cdots$};

  % example n = ExpConvN
  \foreach \k in {-6,-5,-4,-3,-2,-1,0,1,2,3,4,5,6,7,8}{
    \draw[very thick,ThemeGreen] (\k,0) -- (\k,1);
    \fill[ThemeGreen] (\k,1) circle (1.6pt);
    \fill[black] (\k,0) circle (0.6pt);
  }

  \foreach \k in {9,10,11,12}{
    \fill[black] (\k,0) circle (0.75pt);
  }

  \draw[densely dashed,ThemeGreen,very thick] (\ExpConvN,0) -- (\ExpConvN,1.25);
  \node[below] at (\ExpConvN,-0.05) {$n$};
  \node[ThemeGreen,above] at (\ExpConvN,1.27) {edge at \(k=n\)};
\end{scope}

\end{tikzpicture}

\end{OrangeCard}
\end{center}

\end{frame}


% --------------------------------------------------------
\begin{frame}{For \(n<0\), there is no overlap}

\begin{center}
\begin{BlueCard}[
  width=0.95\textwidth,
  halign=center,
  top=12pt,
  bottom=12pt
]
\CardTitle{No common nonzero \(k\)}

\vspace{0.12cm}

\begin{tikzpicture}[x=0.48cm,y=1.05cm,every node/.style={font=\scriptsize}]

\pgfmathtruncatemacro{\LastSample}{\ExpConvMax}

% =========================================================
% Top axis: x[k]
% =========================================================
\draw[-{Latex[length=1.5mm]},thick,black]
  (-6.0,2.20) -- (16.0,2.20) node[right] {$k$};

\node[anchor=east,ThemeGreen,font=\large\bfseries]
  at (-6.25,2.78) {\(x[k]\)};

\node[ThemeGray,font=\large] at (-5.35,2.40) {$\cdots$};
\node[ThemeGray,font=\large] at (15.25,2.40) {$\cdots$};

% zero samples before 0
\foreach \k in {-4,-3,-2,-1}{
  \fill[black] (\k,2.20) circle (0.80pt);
}

% x[k] = alpha^k u[k]
\foreach \k in {0,...,\LastSample}{
  \pgfmathsetmacro{\v}{pow(\ExpConvAlpha,\k)}
  \draw[very thick,ThemeGreen] (\k,2.20) -- (\k,2.20+\v);
  \fill[ThemeGreen] (\k,2.20+\v) circle (1.95pt);
  \fill[black] (\k,2.20) circle (0.65pt);
}

\node[ThemeGreen,above] at (0,3.22) {\(1\)};

% =========================================================
% Bottom axis: h[n-k], example n = -2
% =========================================================
\draw[-{Latex[length=1.5mm]},thick,black]
  (-6.0,-0.55) -- (16.0,-0.55) node[right] {$k$};

\node[anchor=east,ThemeOrange,font=\large\bfseries]
  at (-6.25,0.03) {\(h[n-k]\)};

\node[ThemeGray,font=\large] at (-5.35,0.47) {$\cdots$};

% h[n-k] nonzero for k <= n, illustrated with n=-2
\foreach \k in {-6,-5,-4,-3,-2}{
  \draw[very thick,ThemeOrange] (\k,-0.55) -- (\k,0.45);
  \fill[ThemeOrange] (\k,0.45) circle (1.95pt);
  \fill[black] (\k,-0.55) circle (0.65pt);
}

% zero samples after n
\foreach \k in {-1,0,1,2,3,4,5,6,7,8,9,10,11,12,13,14}{
  \fill[black] (\k,-0.55) circle (0.80pt);
}

% sparse tick labels
\foreach \k in {-5,0,5,10,14}{
  \draw[black] (\k,-0.48) -- (\k,-0.62);
  \node[below] at (\k,-0.67) {$\k$};
}

% boundary and gap
\draw[densely dashed,ThemeOrange,very thick] (-2,-0.55) -- (-2,0.70);
\node[ThemeOrange,above] at (-2,0.75) {\(n<0\)};

% \draw[{Latex[length=1.7mm]}-{Latex[length=1.7mm]},very thick,ThemeRed]
%   (-1.75,0.95) -- (-0.20,0.95);
% \node[ThemeRed,font=\bfseries] at (-0.98,1.18) {gap};

% result
\node[ThemeBlue,font=\Large\bfseries] at (11.2,0.60) {\(y[n]=0\)};

\end{tikzpicture}

% \vspace{0.08cm}

% {\large
% \[
% \text{No overlap}\quad\Rightarrow\quad y[n]=0.
% \]
% }

\end{BlueCard}
\end{center}

\end{frame}







% --------------------------------------------------------
\begin{frame}{For \(n\ge 0\), the overlap is from \(k=0\) to \(k=n\)}

\begin{center}
\begin{GreenCard}[
  width=0.95\textwidth,
  halign=center,
  top=12pt,
  bottom=12pt
]
% \CardTitle{Overlap region}

% \vspace{0.12cm}

\begin{tikzpicture}[x=0.48cm,y=1.05cm,every node/.style={font=\scriptsize}]

\pgfmathtruncatemacro{\LastSample}{\ExpConvMax}
\pgfmathtruncatemacro{\FixedN}{\ExpConvN}

% =========================================================
% Top axis: x[k]
% =========================================================
\draw[-{Latex[length=1.5mm]},thick,black]
  (-5.2,2.20) -- (16.0,2.20) node[right] {$k$};

\node[anchor=east,ThemeGreen,font=\large\bfseries]
  at (-5.45,2.78) {\(x[k]\)};

\node[ThemeGray,font=\large] at (-4.55,2.40) {$\cdots$};
\node[ThemeGray,font=\large] at (15.25,2.40) {$\cdots$};

% zero samples before 0
\foreach \k in {-4,-3,-2,-1}{
  \fill[black] (\k,2.20) circle (0.80pt);
}

% x[k] = alpha^k u[k]
\foreach \k in {0,...,\LastSample}{
  \pgfmathsetmacro{\v}{pow(\ExpConvAlpha,\k)}
  \draw[very thick,ThemeGreen] (\k,2.20) -- (\k,2.20+\v);
  \fill[ThemeGreen] (\k,2.20+\v) circle (1.95pt);
  \fill[black] (\k,2.20) circle (0.65pt);
}

\node[ThemeGreen,above] at (0,3.22) {\(1\)};

% =========================================================
% Bottom axis: h[n-k]
% =========================================================
\draw[-{Latex[length=1.5mm]},thick,black]
  (-5.2,-0.55) -- (16.0,-0.55) node[right] {$k$};

\node[anchor=east,ThemeOrange,font=\large\bfseries]
  at (-5.45,0.03) {\(h[n-k]\)};

\node[ThemeGray,font=\large] at (-4.55,0.47) {$\cdots$};

% h[n-k] nonzero for k <= n
\foreach \k in {-4,-3,-2,-1,0,1,2,3,4,5,6,7,8}{
  \draw[very thick,ThemeOrange] (\k,-0.55) -- (\k,0.45);
  \fill[ThemeOrange] (\k,0.45) circle (1.95pt);
  \fill[black] (\k,-0.55) circle (0.65pt);
}

% zero after n
\foreach \k in {9,10,11,12,13,14}{
  \fill[black] (\k,-0.55) circle (0.80pt);
}

% sparse tick labels
\foreach \k in {-4,0,5,8,10,14}{
  \draw[black] (\k,-0.48) -- (\k,-0.62);
  \node[below] at (\k,-0.67) {$\k$};
}

% overlap guide lines
\foreach \k in {0,1,2,3,4,5,6,7,8}{
  \draw[densely dashed,very thick,ThemeBlue] (\k,2.20) -- (\k,0.45);
}

% overlap bracket
\draw[{Latex[length=1.8mm]}-{Latex[length=1.8mm]},very thick,ThemeBlue]
  (0,0.95) -- (\FixedN,0.95);

\node[ThemeBlue,font=\bfseries] at (4,1.22) {\(0\le k\le n\)};

% edge at n
\draw[densely dashed,ThemeOrange,very thick]
  (\FixedN,-0.55) -- (\FixedN,0.72);

\node[ThemeOrange,above] at (\FixedN,0.77) {\(k=n\)};

\end{tikzpicture}

\vspace{0.08cm}

\(
\text{Overlap}\quad\Rightarrow\quad x[k]h[n-k]=\alpha^k,\qquad 0\le k\le n.
\)

\end{GreenCard}
\end{center}

\end{frame}



























% --------------------------------------------------------
\begin{frame}{Therefore the product is just the overlapping part of \(x[k]\)}

% \begin{OrangeBox}[
%   calloutshadedorange,
%   fontupper=\large,
%   halign=center
% ]{}
% Inside the overlap, \(h[n-k]=1\).  
% So \(x[k]h[n-k]=\alpha^k\) for \(0\le k\le n\), and \(0\) otherwise.
% \end{OrangeBox}

% \vspace{0.10cm}

\begin{center}
\begin{BlueCard}[width=0.90\textwidth,halign=center,top=10pt,bottom=10pt]
\CardTitle{Product sequence for \(n\ge 0\)}

\vspace{0.05cm}

\begin{tikzpicture}[x=0.46cm,y=1.20cm,every node/.style={font=\scriptsize}]

  \draw[-{Latex[length=1.4mm]},thick,black]
    (-3.2,0) -- (15.8,0) node[right] {$k$};
  \draw[-{Latex[length=1.4mm]},thick,black]
    (0,-0.10) -- (0,1.25);

  \node[ThemeGray,font=\large] at (-2.55,0.18) {$\cdots$};
  \node[ThemeGray,font=\large] at (15.10,0.18) {$\cdots$};

  % zero before 0
  \foreach \k in {-2,-1}{
    \fill[black] (\k,0) circle (0.85pt);
  }

  % product alpha^k from 0 to n
  \foreach \k in {0,1,2,3,4,5,6,7,8}{
    \pgfmathsetmacro{\v}{pow(\ExpConvAlpha,\k)}
    \draw[very thick,ThemeBlue] (\k,0) -- (\k,\v);
    \fill[ThemeBlue] (\k,\v) circle (1.9pt);
    \fill[black] (\k,0) circle (0.65pt);
  }

  % zero after n
  \foreach \k in {9,10,11,12,13,14}{
    \fill[black] (\k,0) circle (0.85pt);
  }

  \foreach \k in {0,1,2,3,4,5,6,7,8,9,10,11,12,13,14}{
    \draw[black] (\k,0.025) -- (\k,-0.025);
  }

  \node[below] at (0,-0.04) {$0$};
  \node[below] at (5,-0.04) {$5$};
  \node[below] at (8,-0.04) {$n\)};

  \draw[densely dashed,ThemeOrange,very thick] (\ExpConvN,0) -- (\ExpConvN,0.55);
  \node[ThemeOrange,above] at (\ExpConvN,0.57) {\(n\)};

  \node[ThemeBlue,font=\bfseries] at (5.0,0.80) {\(x[k]h[n-k]\)};

\end{tikzpicture}

\vspace{0.06cm}

{\large
\[
x[k]h[n-k]
=
\begin{cases}
\alpha^k, & 0\le k\le n,\\
0, & \text{otherwise.}
\end{cases}
\]
}

\end{BlueCard}
\end{center}

\end{frame}







% --------------------------------------------------------
\begin{frame}{Now sum the product over all \(k\)}

\begin{GreenBox}[
  calloutshadedgreen,
  fontupper=\large,
  halign=center
]{}
Only the overlap region contributes to the convolution sum.
\end{GreenBox}

\vspace{0.12cm}

\begin{center}
\begin{OrangeCard}[width=0.88\textwidth,halign=center,top=12pt,bottom=12pt]
\CardTitle{For \(n\ge 0\)}

\vspace{0.08cm}

\(
\begin{aligned}
y[n]
&=
\sum_{k=-\infty}^{\infty}x[k]h[n-k]\\[3pt]
&=
\sum_{k=0}^{n}\alpha^k\\[3pt]
&=
1+\alpha+\alpha^2+\cdots+\alpha^n.
\end{aligned}
\)

% \vspace{0.08cm}

% \begin{BlueBox}[
%   calloutshadedblue,
%   halign=center,
%   fontupper=\large
% ]{}
% This is a finite geometric sum.
% \end{BlueBox}

\end{OrangeCard}
\end{center}

\end{frame}







% --------------------------------------------------------
\begin{frame}{The finite geometric sum gives the final answer}

\begin{BlueBox}[
  calloutshadedblue,
  fontupper=\large,
  halign=center
]{}
For \(0<\alpha<1\),
\(
1+\alpha+\alpha^2+\cdots+\alpha^n
=
\frac{1-\alpha^{n+1}}{1-\alpha}.
\)
\end{BlueBox}

\vspace{0.12cm}

\begin{center}
\begin{GreenCard}[width=0.88\textwidth,halign=center,top=12pt,bottom=12pt]
\CardTitle{Final convolution output}

\vspace{0.08cm}

{\Large
\[
y[n]
=
\begin{cases}
\displaystyle \frac{1-\alpha^{n+1}}{1-\alpha}, & n\ge 0,\\[8pt]
0, & n<0.
\end{cases}
\]
}

% \vspace{0.06cm}

% \begin{OrangeBox}[
%   calloutshadedorange,
%   halign=center,
%   fontupper=\Large
% ]{}
% \[
% \displaystyle
% y[n]=\frac{1-\alpha^{n+1}}{1-\alpha}\,u[n].
% \]
% \end{OrangeBox}

\end{GreenCard}
\end{center}

\end{frame}







% --------------------------------------------------------
\begin{frame}{The output rises and approaches a limiting value}

\begin{center}
\begin{BlueCard}[width=0.92\textwidth,halign=center,top=10pt,bottom=10pt]
\CardTitle{Shape of \(y[n]\)}

\vspace{0.05cm}

\begin{tikzpicture}[x=0.38cm,y=0.65cm,every node/.style={font=\scriptsize}]

  \pgfmathsetmacro{\Limit}{1/(1-\ExpConvAlpha)}

  \draw[-{Latex[length=1.5mm]},thick,black]
    (-4.2,0) -- (22.0,0) node[right] {$n$};

  \draw[-{Latex[length=1.5mm]},thick,black]
    (0,-0.15) -- (0,5.30) node[above] {$y[n]$};

  \node[ThemeGray,font=\large] at (-3.55,0.25) {$\cdots$};
  \node[ThemeGray,font=\large] at (21.20,\Limit) {$\cdots$};

  % zero samples for n<0
  \foreach \n in {-3,-2,-1}{
    \fill[black] (\n,0) circle (0.9pt);
  }

  % y[n]=(1-alpha^{n+1})/(1-alpha)
  \foreach \n in {0,...,20}{
    \pgfmathsetmacro{\v}{(1-pow(\ExpConvAlpha,\n+1))/(1-\ExpConvAlpha)}
    \draw[very thick,ThemeBlue] (\n,0) -- (\n,\v);
    \fill[ThemeBlue] (\n,\v) circle (1.75pt);
    \fill[black] (\n,0) circle (0.6pt);
  }

  % limiting line
  \draw[densely dashed,very thick,ThemeOrange] (0,\Limit) -- (20.6,\Limit);
  \node[ThemeOrange,left] at (0,\Limit) {\(\frac{1}{1-\alpha}\)};

  \foreach \n in {-3,-2,-1,0,1,2,3,4,5,6,7,8,9,10,11,12,13,14,15,16,17,18,19,20}{
    \draw[black] (\n,0.025) -- (\n,-0.025);
  }

  \node[below] at (0,-0.05) {$0$};
  \node[below] at (5,-0.05) {$5$};
  \node[below] at (10,-0.05) {$10$};
  \node[below] at (15,-0.05) {$15$};
  \node[below] at (20,-0.05) {$20$};

  \node[ThemeBlue,font=\bfseries] at (-5.5,3.75)
    {\(\displaystyle y[n]=\frac{1-\alpha^{n+1}}{1-\alpha}u[n]\)};

\end{tikzpicture}

% \vspace{0.06cm}

% \begin{OrangeBox}[
%   calloutshadedorange,
%   halign=center,
%   fontupper=\large
% ]{}
% Here the plot uses \(\alpha=\ExpConvAlpha\). As \(n\) increases, more terms are included, and the sum approaches \(\dfrac{1}{1-\alpha}\).
% \end{OrangeBox}

\end{BlueCard}
\end{center}

\end{frame}








% =========================================================
% Controls for Example 2.5
% =========================================================
\def\LeftConvNegMax{8}    % how far left to show
\def\LeftConvPosMax{6}    % how far right to show
\def\LeftConvPosN{5}      % example positive n
\def\LeftConvNegN{-3}     % example negative n





% --------------------------------------------------------
\begin{frame}{Practice (Example 2.5)}

\begin{OrangeBox}[
  calloutshadedorange,
  fontupper=\large,
  halign=center
]{}
Compute
\[
y[n]=x[n]*h[n],
\qquad
x[n]=2^n u[-n],
\qquad
h[n]=u[n].
\]
\end{OrangeBox}

\vspace{0.10cm}

\begin{columns}[c,totalwidth=\textwidth]

% x[n]
\begin{column}{0.49\textwidth}
\begin{GreenCard}[halign=center,top=10pt,bottom=10pt]
\CardTitle{Input \(x[n]=2^n u[-n]\)}

\vspace{0.04cm}

\begin{tikzpicture}[x=0.43cm,y=1.25cm,every node/.style={font=\scriptsize}]

\pgfmathtruncatemacro{\L}{\LeftConvNegMax}
\pgfmathtruncatemacro{\R}{\LeftConvPosMax}

\draw[-{Latex[length=1.4mm]},thick,black]
  (-9.2,0) -- (9.0,0) node[right] {$n$};
\draw[-{Latex[length=1.4mm]},thick,black]
  (0,-0.10) -- (0,1.28);

\node[ThemeGray,font=\large] at (-8.65,0.20) {$\cdots$};
\node[ThemeGray,font=\large] at (8.30,0.20) {$\cdots$};

% nonzero for n <= 0
\foreach \n in {-\L,...,0}{
  \pgfmathsetmacro{\v}{pow(2,\n)}
  \draw[very thick,ThemeGreen] (\n,0) -- (\n,\v);
  \fill[ThemeGreen] (\n,\v) circle (1.8pt);
  \fill[black] (\n,0) circle (0.65pt);
}

% zero for n > 0
\foreach \n in {1,...,\R}{
  \fill[black] (\n,0) circle (0.80pt);
}

\foreach \n in {-8,-6,-4,-2,0,2,4,6,8}{
  \draw[black] (\n,0.025) -- (\n,-0.025);
}

\node[below] at (-4,-0.04) {$-4$};
\node[below] at (0,-0.04) {$0$};
\node[below] at (4,-0.04) {$4$};

\node[left] at (0,1) {$1$};
\node[ThemeGreen,above] at (-3.0,0.40) {\(x[n]\)};

\end{tikzpicture}

\end{GreenCard}
\end{column}

% h[n]
\begin{column}{0.5\textwidth}
\begin{BlueCard}[halign=center,top=10pt,bottom=10pt]
\CardTitle{Impulse response \(h[n]=u[n]\)}

\vspace{0.04cm}

\begin{tikzpicture}[x=0.43cm,y=1.25cm,every node/.style={font=\scriptsize}]

\pgfmathtruncatemacro{\L}{\LeftConvNegMax}
\pgfmathtruncatemacro{\R}{\LeftConvPosMax}

\draw[-{Latex[length=1.4mm]},thick,black]
  (-9.2,0) -- (9.0,0) node[right] {$n$};
\draw[-{Latex[length=1.4mm]},thick,black]
  (0,-0.10) -- (0,1.28);

\node[ThemeGray,font=\large] at (-8.65,0.20) {$\cdots$};
\node[ThemeGray,font=\large] at (8.30,1.08) {$\cdots$};

% zero for n < 0
\foreach \n in {-\L,...,-1}{
  \fill[black] (\n,0) circle (0.80pt);
}

% nonzero for n >= 0
\foreach \n in {0,...,\R}{
  \draw[very thick,ThemeBlue] (\n,0) -- (\n,1);
  \fill[ThemeBlue] (\n,1) circle (1.8pt);
  \fill[black] (\n,0) circle (0.65pt);
}

\foreach \n in {-8,-6,-4,-2,0,2,4,6,8}{
  \draw[black] (\n,0.025) -- (\n,-0.025);
}

\node[below] at (-4,-0.04) {$-4$};
\node[below] at (0,-0.04) {$0$};
\node[below] at (4,-0.04) {$4$};

\node[left] at (0,1) {$1$};
\node[ThemeBlue,above] at (4.2,1.04) {\(h[n]\)};

\end{tikzpicture}

\end{BlueCard}
\end{column}

\end{columns}

\end{frame}







% --------------------------------------------------------
\begin{frame}{First understand how \(h[n-k]\) looks}

\begin{center}
\begin{OrangeCard}[
  width=0.94\textwidth,
  halign=center,
  top=10pt,
  bottom=10pt
]
\CardTitle{Reflect \(h[k]\), then shift it by \(n\)}

\vspace{0.05cm}

\begin{tikzpicture}[x=0.43cm,y=0.70cm,every node/.style={font=\scriptsize}]

\pgfmathtruncatemacro{\L}{\LeftConvNegMax}
\pgfmathtruncatemacro{\R}{\LeftConvPosMax}
\pgfmathtruncatemacro{\Npos}{\LeftConvPosN}
\pgfmathtruncatemacro{\NposPlus}{\Npos+1}

% ---------------- h[k]
\begin{scope}[yshift=2.30cm]
  \node[anchor=east,ThemeBlue,font=\large\bfseries] at (-9.2,0.55) {\(h[k]\)};
  \draw[-{Latex[length=1.3mm]},thick,black] (-8.8,0) -- (9.0,0) node[right] {$k$};

  \node[ThemeGray,font=\large] at (-8.2,0.22) {$\cdots$};
  \node[ThemeGray,font=\large] at (8.3,1.05) {$\cdots$};

  \foreach \k in {-\L,...,-1}{
    \fill[black] (\k,0) circle (0.75pt);
  }

  \foreach \k in {0,...,\R}{
    \draw[very thick,ThemeBlue] (\k,0) -- (\k,1);
    \fill[ThemeBlue] (\k,1) circle (1.65pt);
    \fill[black] (\k,0) circle (0.60pt);
  }

  \node[below] at (0,-0.05) {$0$};
\end{scope}

% ---------------- h[-k]
\begin{scope}[yshift=0.20cm]
  \node[anchor=east,ThemeOrange,font=\large\bfseries] at (-9.2,0.55) {\(h[-k]\)};
  \draw[-{Latex[length=1.3mm]},thick,black] (-8.8,0) -- (9.0,0) node[right] {$k$};

  \node[ThemeGray,font=\large] at (-8.2,1.05) {$\cdots$};
  \node[ThemeGray,font=\large] at (8.3,0.22) {$\cdots$};

  \foreach \k in {-\L,...,0}{
    \draw[very thick,ThemeOrange] (\k,0) -- (\k,1);
    \fill[ThemeOrange] (\k,1) circle (1.65pt);
    \fill[black] (\k,0) circle (0.60pt);
  }

  \foreach \k in {1,...,\R}{
    \fill[black] (\k,0) circle (0.75pt);
  }

  \node[below] at (0,-0.05) {$0$};
\end{scope}

% ---------------- h[n-k]
\begin{scope}[yshift=-1.90cm]
  \node[anchor=east,ThemeGreen,font=\large\bfseries] at (-9.2,0.55) {\(h[n-k]\)};
  \draw[-{Latex[length=1.3mm]},thick,black] (-8.8,0) -- (9.0,0) node[right] {$k$};

  \node[ThemeGray,font=\large] at (-8.2,1.05) {$\cdots$};
  \node[ThemeGray,font=\large] at (8.3,0.22) {$\cdots$};

  \foreach \k in {-\L,...,\Npos}{
    \draw[very thick,ThemeGreen] (\k,0) -- (\k,1);
    \fill[ThemeGreen] (\k,1) circle (1.65pt);
    \fill[black] (\k,0) circle (0.60pt);
  }

  \foreach \k in {\NposPlus,...,\R}{
    \fill[black] (\k,0) circle (0.75pt);
  }

  \draw[densely dashed,ThemeGreen,very thick] (\Npos,0) -- (\Npos,1.25);
  \node[below] at (\Npos,-0.05) {$n$};
  \node[ThemeGreen,above] at (\Npos,1.27) {edge};
\end{scope}

\end{tikzpicture}

\end{OrangeCard}
\end{center}

\end{frame}










% --------------------------------------------------------
\begin{frame}{For \(n<0\), the overlap stops at \(k=n\)}

\begin{center}
\begin{BlueCard}[
  width=0.95\textwidth,
  halign=center,
  top=12pt,
  bottom=12pt
]
% \CardTitle{Overlap for \(n<0\)}

\vspace{0.06cm}

\begin{tikzpicture}[x=0.48cm,y=1.05cm,every node/.style={font=\scriptsize}]

\pgfmathtruncatemacro{\L}{\LeftConvNegMax}
\pgfmathtruncatemacro{\R}{\LeftConvPosMax}
\pgfmathtruncatemacro{\Nneg}{\LeftConvNegN}
\pgfmathtruncatemacro{\NnegPlus}{\Nneg+1}

% =========================================================
% Top axis: x[k]
% =========================================================
\draw[-{Latex[length=1.5mm]},thick,black]
  (-9.2,2.20) -- (9.0,2.20) node[right] {$k$};

\node[anchor=east,ThemeGreen,font=\large\bfseries]
  at (-9.45,2.80) {\(x[k]\)};

\node[ThemeGray,font=\large] at (-8.65,2.42) {$\cdots$};
\node[ThemeGray,font=\large] at (8.30,2.42) {$\cdots$};

% x[k] nonzero for k <= 0
\foreach \k in {-\L,...,0}{
  \pgfmathsetmacro{\v}{pow(2,\k)}
  \draw[very thick,ThemeGreen] (\k,2.20) -- (\k,2.20+\v);
  \fill[ThemeGreen] (\k,2.20+\v) circle (1.95pt);
  \fill[black] (\k,2.20) circle (0.65pt);
}

% x[k] zero for k > 0
\foreach \k in {1,...,\R}{
  \fill[black] (\k,2.20) circle (0.80pt);
}

\node[ThemeGreen,above] at (0,3.22) {\(1\)};

% =========================================================
% Bottom axis: h[n-k], n < 0
% =========================================================
\draw[-{Latex[length=1.5mm]},thick,black]
  (-9.2,-0.55) -- (9.0,-0.55) node[right] {$k$};

\node[anchor=east,ThemeOrange,font=\large\bfseries]
  at (-9.45,0.05) {\(h[n-k]\)};

\node[ThemeGray,font=\large] at (-8.65,0.50) {$\cdots$};

% h[n-k] nonzero for k <= n
\foreach \k in {-\L,...,\Nneg}{
  \draw[very thick,ThemeOrange] (\k,-0.55) -- (\k,0.45);
  \fill[ThemeOrange] (\k,0.45) circle (1.95pt);
  \fill[black] (\k,-0.55) circle (0.65pt);
}

% h[n-k] zero after n
\foreach \k in {\NnegPlus,...,\R}{
  \fill[black] (\k,-0.55) circle (0.80pt);
}

% sparse ticks
\foreach \k in {-8,-4,0,4,8}{
  \draw[black] (\k,-0.48) -- (\k,-0.62);
  \node[below] at (\k,-0.67) {$\k$};
}

% overlap guides: k <= n
\foreach \k in {-\L,...,\Nneg}{
  \draw[densely dashed,very thick,ThemeBlue] (\k,2.20) -- (\k,0.45);
}

% bracket for overlap
\draw[{Latex[length=1.8mm]}-{Latex[length=1.8mm]},very thick,ThemeBlue]
  ({-\L-2},0.95) -- (\Nneg,0.95);

\node[ThemeBlue,font=\bfseries] at (-11,1.32) {overlap: \(k\le n\)};

% edge at negative n
\draw[densely dashed,ThemeOrange,very thick]
  (\Nneg,-0.55) -- (\Nneg,0.72);

\node[ThemeOrange,above] at ({\Nneg+1},0.77) {\(k=n\)};

\end{tikzpicture}

\vspace{0.08cm}

\(
\text{For }n<0,\qquad
x[k]h[n-k]=2^k\quad\text{for }k\le n.
\)

\end{BlueCard}
\end{center}

\end{frame}









% --------------------------------------------------------
\begin{frame}{For \(n\ge 0\), all nonzero samples of \(x[k]\) overlap}

\begin{center}
\begin{GreenCard}[
  width=0.95\textwidth,
  halign=center,
  top=12pt,
  bottom=12pt
]
% \CardTitle{Overlap for \(n\ge 0\)}

\vspace{0.06cm}

\begin{tikzpicture}[x=0.48cm,y=1.05cm,every node/.style={font=\scriptsize}]

\pgfmathtruncatemacro{\L}{\LeftConvNegMax}
\pgfmathtruncatemacro{\R}{\LeftConvPosMax}
\pgfmathtruncatemacro{\Npos}{\LeftConvPosN}
\pgfmathtruncatemacro{\NposPlus}{\Npos+1}

% =========================================================
% Top axis: x[k]
% =========================================================
\draw[-{Latex[length=1.5mm]},thick,black]
  (-9.2,2.20) -- (9.0,2.20) node[right] {$k$};

\node[anchor=east,ThemeGreen,font=\large\bfseries]
  at (-9.45,2.80) {\(x[k]\)};

\node[ThemeGray,font=\large] at (-8.65,2.42) {$\cdots$};
\node[ThemeGray,font=\large] at (8.30,2.42) {$\cdots$};

% x[k] nonzero for k <= 0
\foreach \k in {-\L,...,0}{
  \pgfmathsetmacro{\v}{pow(2,\k)}
  \draw[very thick,ThemeGreen] (\k,2.20) -- (\k,2.20+\v);
  \fill[ThemeGreen] (\k,2.20+\v) circle (1.95pt);
  \fill[black] (\k,2.20) circle (0.65pt);
}

% x[k] zero for k > 0
\foreach \k in {1,...,\R}{
  \fill[black] (\k,2.20) circle (0.80pt);
}

\node[ThemeGreen,above] at (0,3.22) {\(1\)};

% =========================================================
% Bottom axis: h[n-k], n >= 0
% =========================================================
\draw[-{Latex[length=1.5mm]},thick,black]
  (-9.2,-0.55) -- (9.0,-0.55) node[right] {$k$};

\node[anchor=east,ThemeOrange,font=\large\bfseries]
  at (-9.45,0.05) {\(h[n-k]\)};

\node[ThemeGray,font=\large] at (-8.65,0.50) {$\cdots$};

% h[n-k] nonzero for k <= n
\foreach \k in {-\L,...,\Npos}{
  \draw[very thick,ThemeOrange] (\k,-0.55) -- (\k,0.45);
  \fill[ThemeOrange] (\k,0.45) circle (1.95pt);
  \fill[black] (\k,-0.55) circle (0.65pt);
}

% h[n-k] zero after n
\foreach \k in {\NposPlus,...,\R}{
  \fill[black] (\k,-0.55) circle (0.80pt);
}

% sparse ticks
\foreach \k in {-8,-4,0,4,8}{
  \draw[black] (\k,-0.48) -- (\k,-0.62);
  \node[below] at (\k,-0.67) {$\k$};
}

% overlap guides: all k <= 0
\foreach \k in {-\L,...,0}{
  \draw[densely dashed,very thick,ThemeBlue] (\k,2.20) -- (\k,0.45);
}

% bracket for overlap
\draw[{Latex[length=1.8mm]}-{Latex[length=1.8mm]},very thick,ThemeBlue]
  ({-\L-2},0.95) -- (0,0.95);

\node[ThemeBlue,font=\bfseries] at (-11,1.32) {overlap: \(k\le 0\)};

% edge at positive n
\draw[densely dashed,ThemeOrange,very thick]
  (\Npos,-0.55) -- (\Npos,0.72);

\node[ThemeOrange,above] at (\Npos,0.77) {\(k=n\)};

\end{tikzpicture}

\vspace{0.08cm}

\(
\text{For }n\ge 0,\qquad
x[k]h[n-k]=2^k\quad\text{for }k\le 0.
\)

\end{GreenCard}
\end{center}

\end{frame}








% --------------------------------------------------------
\begin{frame}{For \(n<0\), the sum stops at \(k=n\)}

\begin{center}
\begin{BlueCard}[
  width=0.88\textwidth,
  halign=center,
  top=12pt,
  bottom=12pt
]
\CardTitle{Case 1: \(n<0\)}

\vspace{0.08cm}

\[
\begin{aligned}
y[n]
&=
\sum_{k=-\infty}^{\infty}x[k]h[n-k]\\[3pt]
&=
\sum_{k=-\infty}^{n}2^k.
\end{aligned}
\]

\vspace{0.10cm}

\[
\sum_{k=-\infty}^{n}2^k
=
2^n+2^{n-1}+2^{n-2}+\cdots
=
2^{n+1}.
\]

% \vspace{0.08cm}

% \begin{OrangeBox}[
%   calloutshadedorange,
%   halign=center,
%   fontupper=\large
% ]{}
% So for \(n<0\), \(\;y[n]=2^{n+1}\).
% \end{OrangeBox}

\end{BlueCard}
\end{center}

\end{frame}










% --------------------------------------------------------
\begin{frame}{For \(n\ge 0\), the sum includes all left-sided samples}

\begin{center}
\begin{OrangeCard}[
  width=0.88\textwidth,
  halign=center,
  top=12pt,
  bottom=12pt
]
\CardTitle{Case 1: \(n\ge 0\)}

\vspace{0.08cm}

\[
\begin{aligned}
y[n]
&=
\sum_{k=-\infty}^{\infty}x[k]h[n-k]\\[3pt]
&=
\sum_{k=-\infty}^{0}2^k.
\end{aligned}
\]

\vspace{0.10cm}

\[
\sum_{k=-\infty}^{0}2^k
=
\cdots+\frac{1}{8}+\frac{1}{4}+\frac{1}{2}+1
=
2.
\]

% \vspace{0.08cm}

% \begin{BlueBox}[
%   calloutshadedblue,
%   halign=center,
%   fontupper=\large
% ]{}
% So for \(n\ge 0\), \(\;y[n]=2\).
% \end{BlueBox}

\end{OrangeCard}
\end{center}

\end{frame}













% --------------------------------------------------------
\begin{frame}{The output is left-sided exponential, then a constant}

\begin{center}
\begin{GreenCard}[
  width=0.92\textwidth,
  halign=center,
  top=10pt,
  bottom=10pt
]
\CardTitle{Final answer}

\vspace{0.06cm}

\[
y[n]=
\begin{cases}
2^{n+1}, & n<0,\\[4pt]
2, & n\ge 0.
\end{cases}
\]

\vspace{0.06cm}

\begin{tikzpicture}[x=0.48cm,y=0.58cm,every node/.style={font=\scriptsize}]

\pgfmathtruncatemacro{\L}{\LeftConvNegMax}
\pgfmathtruncatemacro{\R}{\LeftConvPosMax}

\draw[-{Latex[length=1.5mm]},thick,black]
  (-9.2,0) -- (9.2,0) node[right] {$n$};

\draw[-{Latex[length=1.5mm]},thick,black]
  (0,-0.15) -- (0,2.65) node[above] {$y[n]$};

\node[ThemeGray,font=\large] at (-8.65,0.25) {$\cdots$};
\node[ThemeGray,font=\large] at (8.50,2.08) {$\cdots$};

% n < 0: 2^{n+1}
\foreach \n in {-\L,...,-1}{
  \pgfmathsetmacro{\v}{pow(2,\n+1)}
  \draw[very thick,ThemeGreen] (\n,0) -- (\n,\v);
  \fill[ThemeGreen] (\n,\v) circle (1.9pt);
  \fill[black] (\n,0) circle (0.65pt);
}

% n >= 0: constant 2
\foreach \n in {0,...,\R}{
  \draw[very thick,ThemeGreen] (\n,0) -- (\n,2);
  \fill[ThemeGreen] (\n,2) circle (1.9pt);
  \fill[black] (\n,0) circle (0.65pt);
}

% limiting/constant line
\draw[densely dashed,very thick,ThemeOrange] (0,2) -- (8.3,2);
\node[ThemeOrange,left] at (0,2) {$2$};

% ticks
\foreach \n in {-8,-4,0,4,8}{
  \draw[black] (\n,0.035) -- (\n,-0.035);
  \node[below] at (\n,-0.06) {$\n$};
}

\node[left] at (0,1) {$1$};
\node[ThemeGreen,font=\bfseries] at (4.7,2.35) {\(y[n]\)};

\end{tikzpicture}

\end{GreenCard}
\end{center}

\end{frame}





% =========================================================
% Controls for continuous-time convolution example
% =========================================================
\newcommand{\CTConvA}{0.70}       % plotted value of a
\newcommand{\CTConvTpos}{3.00}    % example positive t
\newcommand{\CTConvTneg}{-1.25}   % example negative t



% --------------------------------------------------------
\begin{frame}{Practice (Example 2.6)}

\begin{OrangeBox}[
  calloutshadedorange,
  fontupper=\large,
  halign=center
]{}
Let \(a>0\). Compute
\[
y(t)=x(t)*h(t),
\qquad
x(t)=e^{-at}u(t),
\qquad
h(t)=u(t).
\]
\end{OrangeBox}

\vspace{0.02cm}

\begin{columns}[c,totalwidth=\textwidth]

% x(t)
\begin{column}{0.49\textwidth}
\begin{GreenCard}[halign=center,top=10pt,bottom=10pt]
\CardTitle{Input \(x(t)=e^{-at}u(t)\)}

\vspace{0.04cm}

\begin{tikzpicture}[x=0.78cm,y=1.25cm,every node/.style={font=\scriptsize}]

  \draw[-{Latex[length=1.5mm]},thick,black]
    (-2.4,0) -- (5.1,0) node[right] {$t$};
  \draw[-{Latex[length=1.5mm]},thick,black]
    (0,-0.10) -- (0,1.25);

  % zero for t<0
  \draw[very thick,ThemeGray] (-2.0,0) -- (0,0);

  % exponential for t>=0
  \draw[very thick,ThemeGreen,smooth]
    plot[domain=0:4.6,samples=120] (\x,{exp(-\CTConvA*\x)});

  \draw[very thick,ThemeGreen] (0,0) -- (0,1);

  \node[ThemeGray,font=\large] at (-2.05,0.18) {$\cdots$};
  \node[ThemeGray,font=\large] at (4.65,0.18) {$\cdots$};

  \node[below] at (0,-0.04) {$0$};
  \node[left] at (0,1) {$1$};

  \node[ThemeGreen,above] at (1.75,0.45) {\(e^{-at}u(t)\)};

\end{tikzpicture}

% \vspace{0.04cm}

% {\normalsize plotted with \(a=\CTConvA\)}

\end{GreenCard}
\end{column}

% h(t)
\begin{column}{0.49\textwidth}
\begin{BlueCard}[halign=center,top=10pt,bottom=10pt]
\CardTitle{Impulse response \(h(t)=u(t)\)}

\vspace{0.04cm}

\begin{tikzpicture}[x=0.78cm,y=1.25cm,every node/.style={font=\scriptsize}]

  \draw[-{Latex[length=1.5mm]},thick,black]
    (-2.4,0) -- (5.1,0) node[right] {$t$};
  \draw[-{Latex[length=1.5mm]},thick,black]
    (0,-0.10) -- (0,1.25);

  % unit step
  \draw[very thick,ThemeGray] (-2.0,0) -- (0,0);
  \draw[very thick,ThemeBlue] (0,0) -- (0,1);
  \draw[very thick,ThemeBlue] (0,1) -- (4.6,1);

  \node[ThemeGray,font=\large] at (-2.05,0.18) {$\cdots$};
  \node[ThemeGray,font=\large] at (4.65,1.08) {$\cdots$};

  \node[below] at (0,-0.04) {$0$};
  \node[left] at (0,1) {$1$};

  \node[ThemeBlue,above] at (2.55,1.04) {\(u(t)\)};

\end{tikzpicture}

% \vspace{0.04cm}

% {\normalsize a unit step}

\end{BlueCard}
\end{column}

\end{columns}

\end{frame}







% --------------------------------------------------------
\begin{frame}{First understand how \(h(t-\tau)\) moves}

\begin{center}
\begin{OrangeCard}[
  width=0.94\textwidth,
  halign=center,
  top=10pt,
  bottom=10pt
]
\CardTitle{Reflect \(h(\tau)\), then shift by \(t\)}

\vspace{0.06cm}

\begin{tikzpicture}[x=0.83cm,y=0.74cm,every node/.style={font=\scriptsize}]

% =========================================================
% h(tau)
% =========================================================
\begin{scope}[yshift=2.35cm]
  \node[anchor=east,ThemeBlue,font=\large\bfseries] at (-3.0,0.55) {\(h(\tau)\)};

  \draw[-{Latex[length=1.4mm]},thick,black]
    (-2.7,0) -- (5.1,0) node[right] {$\tau$};

  \draw[very thick,ThemeGray] (-2.3,0) -- (0,0);
  \draw[very thick,ThemeBlue] (0,0) -- (0,1);
  \draw[very thick,ThemeBlue] (0,1) -- (4.6,1);

  \node[ThemeGray,font=\large] at (-2.25,0.20) {$\cdots$};
  \node[ThemeGray,font=\large] at (4.60,1.10) {$\cdots$};

  \node[below] at (0,-0.05) {$0$};
\end{scope}

% =========================================================
% h(-tau)
% =========================================================
\begin{scope}[yshift=0.30cm]
  \node[anchor=east,ThemeOrange,font=\large\bfseries] at (-3.0,0.55) {\(h(-\tau)\)};

  \draw[-{Latex[length=1.4mm]},thick,black]
    (-5.1,0) -- (2.7,0) node[right] {$\tau$};

  \draw[very thick,ThemeOrange] (-4.6,1) -- (0,1);
  \draw[very thick,ThemeOrange] (0,1) -- (0,0);
  \draw[very thick,ThemeGray] (0,0) -- (2.3,0);

  \node[ThemeGray,font=\large] at (-4.60,1.10) {$\cdots$};
  \node[ThemeGray,font=\large] at (2.25,0.20) {$\cdots$};

  \node[below] at (0,-0.05) {$0$};
\end{scope}

% =========================================================
% h(t-tau)
% =========================================================
\begin{scope}[yshift=-1.75cm]
  \node[anchor=east,ThemeGreen,font=\large\bfseries] at (-3.0,0.55) {\(h(t-\tau)\)};

  \draw[-{Latex[length=1.4mm]},thick,black]
    (-3.0,0) -- (5.3,0) node[right] {$\tau$};

  \draw[very thick,ThemeGreen] (-2.6,1) -- (\CTConvTpos,1);
  \draw[very thick,ThemeGreen] (\CTConvTpos,1) -- (\CTConvTpos,0);
  \draw[very thick,ThemeGray] (\CTConvTpos,0) -- (4.8,0);

  \node[ThemeGray,font=\large] at (-2.55,1.10) {$\cdots$};
  \node[ThemeGray,font=\large] at (4.75,0.20) {$\cdots$};

  \draw[densely dashed,ThemeGreen,very thick]
    (\CTConvTpos,0) -- (\CTConvTpos,1.25);

  \node[below] at (0,-0.05) {$0$};
  \node[below] at (\CTConvTpos,-0.05) {$t$};
  \node[ThemeGreen,above] at (\CTConvTpos,1.27) {edge};

\end{scope}

\end{tikzpicture}

\end{OrangeCard}
\end{center}

\end{frame}







% --------------------------------------------------------
\begin{frame}{For \(t<0\), there is no overlap}

\begin{center}
\begin{BlueCard}[
  width=0.95\textwidth,
  halign=center,
  top=12pt,
  bottom=12pt
]
\CardTitle{No common nonzero \(\tau\)}

\vspace{0.12cm}

\begin{tikzpicture}[x=0.86cm,y=1.05cm,every node/.style={font=\scriptsize}]

% =========================================================
% Top axis: x(tau)
% =========================================================
\draw[-{Latex[length=1.5mm]},thick,black]
  (-2.8,2.20) -- (5.2,2.20) node[right] {$\tau$};

\node[anchor=east,ThemeGreen,font=\large\bfseries]
  at (-3.00,2.78) {\(x(\tau)\)};

\node[ThemeGray,font=\large] at (-2.45,2.38) {$\cdots$};
\node[ThemeGray,font=\large] at (4.75,2.38) {$\cdots$};

\draw[very thick,ThemeGray] (-2.3,2.20) -- (0,2.20);
\draw[very thick,ThemeGreen] (0,2.20) -- (0,3.20);
\draw[very thick,ThemeGreen,smooth]
  plot[domain=0:4.6,samples=120] (\x,{2.20+exp(-\CTConvA*\x)});

\node[below] at (0,2.15) {$0$};
\node[ThemeGreen,above] at (0,3.22) {\(1\)};

% =========================================================
% Bottom axis: h(t-tau), t<0
% =========================================================
\draw[-{Latex[length=1.5mm]},thick,black]
  (-2.8,-0.55) -- (5.2,-0.55) node[right] {$\tau$};

\node[anchor=east,ThemeOrange,font=\large\bfseries]
  at (-3.00,0.03) {\(h(t-\tau)\)};

\node[ThemeGray,font=\large] at (-2.45,0.50) {$\cdots$};

\draw[very thick,ThemeOrange] (-2.3,0.45) -- (\CTConvTneg,0.45);
\draw[very thick,ThemeOrange] (\CTConvTneg,0.45) -- (\CTConvTneg,-0.55);
\draw[very thick,ThemeGray] (\CTConvTneg,-0.55) -- (4.7,-0.55);

\draw[densely dashed,ThemeOrange,very thick]
  (\CTConvTneg,-0.55) -- (\CTConvTneg,0.75);

\node[below] at (\CTConvTneg,-0.60) {$t<0$};
\node[below] at (0,-0.60) {$0$};

% % gap
% \draw[{Latex[length=1.7mm]}-{Latex[length=1.7mm]},very thick,ThemeRed]
%   ({\CTConvTneg+0.20},0.95) -- (-0.20,0.95);

% \node[ThemeRed,font=\bfseries] at ({(\CTConvTneg-0.20)/2},1.20) {gap};

% result
\node[ThemeBlue,font=\Large\bfseries] at (3.45,0.60) {\(y(t)=0\)};

\end{tikzpicture}

\vspace{0.08cm}

\(
\text{No overlap}\quad\Rightarrow\quad y(t)=0.
\)

\end{BlueCard}
\end{center}

\end{frame}









% --------------------------------------------------------
\begin{frame}{For \(t>0\), the overlap is \(0\le \tau\le t\)}

\begin{center}
\begin{GreenCard}[
  width=0.95\textwidth,
  halign=center,
  top=12pt,
  bottom=12pt
]
% \CardTitle{Overlap region}

% \vspace{0.12cm}

\begin{tikzpicture}[x=0.86cm,y=1.05cm,every node/.style={font=\scriptsize}]

% =========================================================
% Top axis: x(tau)
% =========================================================
\draw[-{Latex[length=1.5mm]},thick,black]
  (-2.4,2.20) -- (5.4,2.20) node[right] {$\tau$};

\node[anchor=east,ThemeGreen,font=\large\bfseries]
  at (-2.65,2.78) {\(x(\tau)\)};

\node[ThemeGray,font=\large] at (-2.05,2.38) {$\cdots$};
\node[ThemeGray,font=\large] at (4.95,2.38) {$\cdots$};

% zero part of x(tau), then jump, then decay
\draw[very thick,ThemeGray] (-2.0,2.20) -- (0,2.20);
\draw[very thick,ThemeGreen] (0,2.20) -- (0,3.20);
\draw[very thick,ThemeGreen,smooth]
  plot[domain=0:4.8,samples=120] (\x,{2.20+exp(-\CTConvA*\x)});

\node[below] at (0,2.15) {$0$};
\node[ThemeGreen,above] at (0,3.22) {\(1\)};

% =========================================================
% Bottom axis: h(t-tau), t>0
% =========================================================
\draw[-{Latex[length=1.5mm]},thick,black]
  (-2.4,-0.55) -- (5.4,-0.55) node[right] {$\tau$};

\node[anchor=east,ThemeOrange,font=\large\bfseries]
  at (-2.65,0.03) {\(h(t-\tau)\)};

\node[ThemeGray,font=\large] at (-2.05,0.50) {$\cdots$};

% h(t-tau): equal to 1 for tau <= t, zero after t
\draw[very thick,ThemeOrange] (-2.0,0.45) -- (\CTConvTpos,0.45);
\draw[very thick,ThemeOrange] (\CTConvTpos,0.45) -- (\CTConvTpos,-0.55);
\draw[very thick,ThemeGray] (\CTConvTpos,-0.55) -- (4.9,-0.55);

\node[below] at (0,-0.60) {$0$};
\node[below] at (\CTConvTpos,-0.60) {$t$};

% vertical guides showing overlap interval
\draw[densely dashed,very thick,ThemeBlue]
  (0,2.20) -- (0,0.45);

\draw[densely dashed,very thick,ThemeBlue]
  (\CTConvTpos,2.20) -- (\CTConvTpos,0.45);

% overlap bracket
\draw[{Latex[length=1.8mm]}-{Latex[length=1.8mm]},very thick,ThemeBlue]
  (0,0.95) -- (\CTConvTpos,0.95);

\node[ThemeBlue,font=\bfseries] at ({0.5*\CTConvTpos},1.22)
  {\(0\le \tau\le t\)};

\end{tikzpicture}

\vspace{0.08cm}

\(
\text{Overlap}\quad\Rightarrow\quad
x(\tau)h(t-\tau)=e^{-a\tau},
\qquad 0\le \tau\le t.
\) 

\end{GreenCard}
\end{center}

\end{frame}













% --------------------------------------------------------
\begin{frame}{Now integrate only over the overlap}

\begin{center}
\begin{OrangeCard}[
  width=0.88\textwidth,
  halign=center,
  top=12pt,
  bottom=12pt
]
% \CardTitle{For \(t>0\)}

% \vspace{0.08cm}

\[
\begin{aligned}
\text{For } t>0, \; y(t)
&=
\int_{-\infty}^{\infty}x(\tau)h(t-\tau)\,d\tau\\[3pt]
&=
\int_{0}^{t} e^{-a\tau}\,d\tau\\[3pt]
&=
\left[-\frac{1}{a}e^{-a\tau}\right]_{0}^{t}\\[3pt]
&=
\frac{1}{a}\left(1-e^{-at}\right).
\end{aligned}
\]

\vspace{0.08cm}

\begin{BlueBox}[
  calloutshadedblue,
  halign=center,
  fontupper=\large
]{}
For \(t<0\), there is no overlap, so \(y(t)=0\).
\end{BlueBox}

\end{OrangeCard}
\end{center}

\end{frame}






% --------------------------------------------------------
\begin{frame}{The output gradually rises to \(1/a\)}

\begin{center}
\begin{GreenCard}[
  width=\textwidth,
  halign=center,
  top=10pt,
  bottom=10pt
]
% \CardTitle{Final answer}

% \vspace{0.06cm}

\[
y(t)
=
\frac{1}{a}\left(1-e^{-at}\right)u(t).
\]

\vspace{0.06cm}

\begin{tikzpicture}[x=0.90cm,y=1.15cm,every node/.style={font=\scriptsize}]

\draw[-{Latex[length=1.5mm]},thick,black]
  (-2.3,0) -- (6.0,0) node[right] {$t$};

\draw[-{Latex[length=1.5mm]},thick,black]
  (0,-0.12) -- (0,1.75) node[above] {$y(t)$};

% visible zero part for t < 0
\draw[very thick,ThemeGreen] (-2.0,0) -- (0,0);

% output for t >= 0
\draw[very thick,ThemeGreen,smooth]
  plot[domain=0:5.4,samples=140] (\x,{1.35*(1-exp(-\CTConvA*\x))});

% limiting value
\draw[densely dashed,very thick,ThemeOrange]
  (0,1.35) -- (5.5,1.35);

\node[ThemeOrange,left] at (0,1.35) {\(\frac{1}{a}\)};
\node[below] at (0,-0.05) {$0$};

\node[ThemeGray,font=\large] at (-1.85,0.18) {$\cdots$};
\node[ThemeGray,font=\large] at (5.55,1.35) {$\cdots$};

\node[ThemeGreen,font=\bfseries] at (6.20,0.72)
  {\(\displaystyle y(t)=\frac{1}{a}(1-e^{-at})u(t)\)};

\end{tikzpicture}

\vspace{0.08cm}

\begin{OrangeBox}[
  calloutshadedorange,
  halign=center,
  fontupper=\large
]{}
As \(t\) increases, the overlap grows, so the integral rises toward \(\dfrac{1}{a}\).
\end{OrangeBox}

\end{GreenCard}
\end{center}

\end{frame}






% --------------------------------------------------------
\begin{frame}{Property: Convolution distributes over addition}

\begin{BlueBox}[
  calloutshadedblue,
  fontupper=\large,
  halign=center
]{}
If one signal is a sum of two simpler signals, \\ we may convolve the two parts separately.
\end{BlueBox}

\vspace{0.14cm}

\begin{center}
\begin{GreenCard}[
  width=0.88\textwidth,
  halign=center,
  top=12pt,
  bottom=12pt
]
% \CardTitle{Distributive property}

% \vspace{0.06cm}

{\Large
\[
x[n]*\big(h_1[n]+h_2[n]\big)
=
x[n]*h_1[n]+x[n]*h_2[n].
\]
}

\vspace{0.08cm}

{\Large
\[
\big(x_1[n]+x_2[n]\big)*h[n]
=
x_1[n]*h[n]+x_2[n]*h[n].
\]
}

\end{GreenCard}
\end{center}

% \vspace{0.12cm}

% \begin{OrangeBox}[
%   calloutshadedorange,
%   fontupper=\large,
%   halign=center
% ]{}
% This is useful when one signal looks complicated, but can be split into easier pieces.
% \end{OrangeBox}

\end{frame}








% --------------------------------------------------------
\begin{frame}{Practice (Example 2.10)}

\begin{OrangeBox}[
  calloutshadedorange,
  fontupper=\large,
  halign=center
]{}
Compute
\[
y[n]=x[n]*h[n],
\qquad
x[n]=\left(\frac12\right)^n u[n]+2^n u[-n],
\qquad
h[n]=u[n].
\]
\end{OrangeBox}

\vspace{0.10cm}

\begin{columns}[c,totalwidth=\textwidth]

% x[n]
\begin{column}{0.58\textwidth}
\begin{GreenCard}[halign=center,top=10pt,bottom=10pt]
\CardTitle{Input \(x[n]\)}

\vspace{0.04cm}

\begin{tikzpicture}[x=0.43cm,y=0.58cm,every node/.style={font=\scriptsize}]

\draw[-{Latex[length=1.4mm]},thick,black]
  (-7.2,0) -- (8.0,0) node[right] {$n$};
\draw[-{Latex[length=1.4mm]},thick,black]
  (0,-0.15) -- (0,2.65);

\node[ThemeGray,font=\large] at (-6.75,0.25) {$\cdots$};
\node[ThemeGray,font=\large] at (7.35,0.25) {$\cdots$};

% left part: 2^n u[-n]
\foreach \n in {-6,-5,-4,-3,-2,-1}{
  \pgfmathsetmacro{\v}{pow(2,\n)}
  \draw[very thick,ThemeOrange] (\n,0) -- (\n,\v);
  \fill[ThemeOrange] (\n,\v) circle (1.8pt);
  \fill[black] (\n,0) circle (0.65pt);
}

% n=0 has both contributions, so x[0]=2
\draw[very thick,ThemeGreen] (0,0) -- (0,2);
\fill[ThemeGreen] (0,2) circle (1.9pt);
\fill[black] (0,0) circle (0.65pt);

% right part: (1/2)^n u[n]
\foreach \n in {1,2,3,4,5,6,7}{
  \pgfmathsetmacro{\v}{pow(0.5,\n)}
  \draw[very thick,ThemeBlue] (\n,0) -- (\n,\v);
  \fill[ThemeBlue] (\n,\v) circle (1.8pt);
  \fill[black] (\n,0) circle (0.65pt);
}

\foreach \n in {-6,-3,0,3,6}{
  \draw[black] (\n,0.04) -- (\n,-0.04);
  \node[below] at (\n,-0.08) {$\n$};
}

\node[left] at (0,1) {$1$};
\node[left] at (0,2) {$2$};

\node[ThemeOrange,font=\bfseries] at (-3.8,0.72) {\(2^n u[-n]\)};
\node[ThemeBlue,font=\bfseries] at (4.2,0.72) {\(\left(\frac12\right)^n u[n]\)};

\end{tikzpicture}

\end{GreenCard}
\end{column}

% h[n]
\begin{column}{0.38\textwidth}
\begin{BlueCard}[halign=center,top=10pt,bottom=10pt]
\CardTitle{\(h[n]=u[n]\)}

\vspace{0.04cm}

\begin{tikzpicture}[x=0.45cm,y=0.70cm,every node/.style={font=\scriptsize}]

\draw[-{Latex[length=1.4mm]},thick,black]
  (-4.2,0) -- (6.5,0) node[right] {$n$};
\draw[-{Latex[length=1.4mm]},thick,black]
  (0,-0.12) -- (0,1.45);

\node[ThemeGray,font=\large] at (-3.8,0.20) {$\cdots$};
\node[ThemeGray,font=\large] at (6.0,1.08) {$\cdots$};

\foreach \n in {-3,-2,-1}{
  \fill[black] (\n,0) circle (0.75pt);
}

\foreach \n in {0,1,2,3,4,5}{
  \draw[very thick,ThemeBlue] (\n,0) -- (\n,1);
  \fill[ThemeBlue] (\n,1) circle (1.7pt);
  \fill[black] (\n,0) circle (0.65pt);
}

\node[below] at (0,-0.06) {$0$};
\node[left] at (0,1) {$1$};

\end{tikzpicture}

\end{BlueCard}
\end{column}

\end{columns}

\end{frame}





% --------------------------------------------------------
\begin{frame}{Oops, different input expressions on the left and right side!}

\begin{center}
\begin{OrangeCard}[
  width=0.94\textwidth,
  halign=center,
  top=12pt,
  bottom=12pt
]
% \CardTitle{The issue}

% \vspace{0.08cm}

\begin{tikzpicture}[x=0.52cm,y=0.66cm,every node/.style={font=\scriptsize}]

% axis
\draw[-{Latex[length=1.4mm]},thick,black]
  (-7.0,0) -- (8.0,0) node[right] {$k$};
\draw[-{Latex[length=1.4mm]},thick,black]
  (0,-0.12) -- (0,2.45);

\node[ThemeGray,font=\large] at (-6.5,0.25) {$\cdots$};
\node[ThemeGray,font=\large] at (7.35,0.25) {$\cdots$};

% left part x[k]
\foreach \k in {-6,-5,-4,-3,-2,-1}{
  \pgfmathsetmacro{\v}{pow(2,\k)}
  \draw[very thick,ThemeOrange] (\k,0) -- (\k,\v);
  \fill[ThemeOrange] (\k,\v) circle (1.7pt);
  \fill[black] (\k,0) circle (0.65pt);
}

% at zero
\draw[very thick,ThemeGreen] (0,0) -- (0,2);
\fill[ThemeGreen] (0,2) circle (1.8pt);
\fill[black] (0,0) circle (0.65pt);

% right part x[k]
\foreach \k in {1,2,3,4,5,6,7}{
  \pgfmathsetmacro{\v}{pow(0.5,\k)}
  \draw[very thick,ThemeBlue] (\k,0) -- (\k,\v);
  \fill[ThemeBlue] (\k,\v) circle (1.7pt);
  \fill[black] (\k,0) circle (0.65pt);
}

% split line
\draw[densely dashed,very thick,ThemeRed] (0,0) -- (0,2.25);
\node[ThemeRed,above] at (0,2.28) {formula changes};

\node[ThemeOrange,font=\bfseries] at (-3.6,0.70) {\(2^k\)};
\node[ThemeBlue,font=\bfseries] at (3.9,0.70) {\(\left(\frac12\right)^k\)};

\node[below] at (0,-0.08) {$0$};

\end{tikzpicture}

\vspace{0.10cm}

\(
x[k]=
\underbrace{2^k u[-k]}_{\text{left-sided part}}
+
\underbrace{\left(\frac12\right)^k u[k]}_{\text{right-sided part}}.
\)

\vspace{0.04cm}

\begin{BlueBox}[
  calloutshadedblue,
  halign=center,
  fontupper=\large
]{}
So, we split it first.
\end{BlueBox}

\end{OrangeCard}
\end{center}

\end{frame}







% --------------------------------------------------------
\begin{frame}{Split the problem into two easier ones!}

\begin{center}
\begin{GreenCard}[
  width=\textwidth,
  halign=center,
  top=12pt,
  bottom=12pt
]
\CardTitle{Split \(x[n]\), then convolve separately}

\vspace{0.08cm}

\[
x[n]=x_1[n]+x_2[n],
\qquad
x_1[n]=\left(\frac12\right)^n u[n],
\qquad
x_2[n]=2^n u[-n].
\]

\vspace{0.12cm}

\(
\begin{aligned}
y[n]
&=x[n]*h[n]\\[3pt]
% &=\big(x_1[n]+x_2[n]\big)*h[n]\\[3pt]
&=x_1[n]*h[n]+x_2[n]*h[n]\\[3pt]
&=y_1[n]+y_2[n].
\end{aligned}
\)

\vspace{0.08cm}

\begin{OrangeBox}[
  calloutshadedorange,
  halign=center,
  fontupper=\large
]{}
Each part is one we already know how to do.
\end{OrangeBox}

\end{GreenCard}
\end{center}

\end{frame}




% --------------------------------------------------------
\begin{frame}{The two smaller convolutions have simple answers}

\begin{columns}[c,totalwidth=\textwidth]






% y2
\begin{column}{0.49\textwidth}
\begin{OrangeCard}[halign=center,top=10pt,bottom=10pt]
\CardTitle{Left-sided part}

\vspace{0.04cm}

\[
x_2[n]=2^n u[-n],
\qquad
h[n]=u[n].
\]

\vspace{0.06cm}

\[
y_2[n]=x_2[n]*h[n].
\]

\vspace{0.06cm}

\[
y_2[n]=
\begin{cases}
2^{n+1}, & n<0,\\[3pt]
2, & n\ge 0.
\end{cases}
\]

% \vspace{0.04cm}

% {\normalsize
% This is the accumulated sum of a left-sided exponential.
% }

\end{OrangeCard}
\end{column}






% y1
\begin{column}{0.49\textwidth}
\begin{BlueCard}[halign=center,top=10pt,bottom=10pt]
\CardTitle{Right-sided part}

\vspace{0.04cm}

\[
x_1[n]=\left(\frac12\right)^n u[n],
\qquad
h[n]=u[n].
\]

\vspace{0.06cm}

\[
y_1[n]=x_1[n]*h[n].
\]

\vspace{0.06cm}

\[
y_1[n]
=
\left(2-2^{-n}\right)u[n].
\]

% \vspace{0.04cm}

% {\normalsize
% This is the accumulated sum of a decaying right-sided exponential.
% }

\end{BlueCard}
\end{column}


\end{columns}

\end{frame}








% --------------------------------------------------------
\begin{frame}{Now add the two outputs}

\begin{center}
\begin{GreenCard}[
  width=0.90\textwidth,
  halign=center,
  top=12pt,
  bottom=12pt
]
\CardTitle{Final expression}

\vspace{0.06cm}

\[
y[n]=y_1[n]+y_2[n].
\]

\vspace{0.08cm}

% \[
% y[n]
% =
% \begin{cases}
% 2^{n+1}, & n<0,\\[5pt]
% \left(2-2^{-n}\right)+2, & n\ge 0.
% \end{cases}
% \]

% \vspace{0.08cm}

\[
\boxed{
y[n]
=
\begin{cases}
2^{n+1}, & n<0,\\[5pt]
4-2^{-n}, & n\ge 0.
\end{cases}}
\]

\vspace{0.06cm}

\begin{OrangeBox}[
  calloutshadedorange,
  halign=center,
  fontupper=\large
]{}
The left side rises toward \(1\), then the right side rises toward \(4\).
\end{OrangeBox}

\end{GreenCard}
\end{center}

\end{frame}





% --------------------------------------------------------
\begin{frame}{The output approaches \(4\) on the right}

\begin{center}
\begin{BlueCard}[
  width=0.92\textwidth,
  halign=center,
  top=10pt,
  bottom=10pt
]
\CardTitle{Sketch of \(y[n]\)}

\vspace{0.04cm}

\begin{tikzpicture}[x=0.55cm,y=0.55cm,every node/.style={font=\scriptsize}]

\draw[-{Latex[length=1.5mm]},thick,black]
  (-7.2,0) -- (9.0,0) node[right] {$n$};

\draw[-{Latex[length=1.5mm]},thick,black]
  (0,-0.15) -- (0,4.55) node[above] {$y[n]$};

\node[ThemeGray,font=\large] at (-6.7,0.25) {$\cdots$};
\node[ThemeGray,font=\large] at (8.35,4.05) {$\cdots$};

% n < 0: 2^{n+1}
\foreach \n in {-6,-5,-4,-3,-2,-1}{
  \pgfmathsetmacro{\v}{pow(2,\n+1)}
  \draw[very thick,ThemeGreen] (\n,0) -- (\n,\v);
  \fill[ThemeGreen] (\n,\v) circle (1.9pt);
  \fill[black] (\n,0) circle (0.65pt);
}

% n >= 0: 4 - 2^{-n}
\foreach \n in {0,1,2,3,4,5,6,7}{
  \pgfmathsetmacro{\v}{4-pow(2,-\n)}
  \draw[very thick,ThemeGreen] (\n,0) -- (\n,\v);
  \fill[ThemeGreen] (\n,\v) circle (1.9pt);
  \fill[black] (\n,0) circle (0.65pt);
}

% asymptote
\draw[densely dashed,very thick,ThemeOrange]
  (0,4) -- (8.0,4);

\node[ThemeOrange,left] at (0,4) {$4$};

% ticks
\foreach \n in {-6,-3,0,3,6}{
  \draw[black] (\n,0.04) -- (\n,-0.04);
  \node[below] at (\n,-0.08) {$\n$};
}

\node[left] at (0,1) {$1$};
\node[left] at (0,2) {$2$};
\node[left] at (0,3) {$3$};

\node[ThemeGreen,font=\bfseries] at (4.4,3.15)
  {\(4-2^{-n}\)};

\node[ThemeGreen,font=\bfseries] at (-3.6,0.95)
  {\(2^{n+1}\)};

\end{tikzpicture}

\vspace{0.08cm}

\[
y[n]
=
\begin{cases}
2^{n+1}, & n<0,\\[4pt]
4-2^{-n}, & n\ge 0.
\end{cases}
\]

\end{BlueCard}
\end{center}

\end{frame}





% --------------------------------------------------------
\begin{frame}[plain]

\vfill

\begin{center}

{\fontsize{30}{34}\selectfont\bfseries\color{FrameTitleColor}
Causality and Stability\par}

% \vspace{0.25cm}

% {\Large\color{ThemeGray}
% Convolution is simply a mathematical operation that appears very universally.\par}

\end{center}

\vfill

\end{frame}





\begin{frame}{Causality: No Time Travel Please!!}

\begin{center}

\begin{figure}[t]
    \centering
    \includegraphics[width=0.9\textwidth]{figs/time-travel.png}
\end{figure}

\end{center}

\end{frame}




% --------------------------------------------------------
\begin{frame}{For an LTI system, causality is visible from \(h(t)\)}

\begin{BlueBox}[
  calloutshadedblue,
  fontupper=\large,
  halign=center
]{}
The impulse response \(h(t)\) shows what the system does after we poke it with \(\delta(t)\) at \(t=0\).
\end{BlueBox}

\vspace{0.12cm}

\begin{columns}[c,totalwidth=\textwidth]

% causal
\begin{column}{0.49\textwidth}
\begin{GreenCard}[halign=center,top=10pt,bottom=10pt]
\CardTitle{Causal}

\vspace{0.04cm}

\begin{tikzpicture}[x=0.78cm,y=1.10cm,every node/.style={font=\scriptsize}]

  \draw[-{Latex[length=1.5mm]},thick,black]
    (-2.6,0) -- (4.0,0) node[right] {$t$};
  \draw[-{Latex[length=1.5mm]},thick,black]
    (0,-0.12) -- (0,1.35);

  % zero before impulse
  \draw[very thick,ThemeGreen] (-2.2,0) -- (0,0);

  % response after 0
  \draw[very thick,ThemeGreen,smooth]
    plot[domain=0:3.4,samples=100] (\x,{exp(-0.75*\x)});

  \draw[densely dashed,ThemeGray] (0,0) -- (0,1.15);

  \node[below] at (0,-0.05) {$0$};
  \node[ThemeGreen,above] at (1.45,0.60) {\(h(t)\)};
  \node[ThemeGray] at (-1.25,0.22) {nothing before};

\end{tikzpicture}

\vspace{0.06cm}

\(
h(t)=0\qquad \text{for }t<0.
\)

\end{GreenCard}
\end{column}

% noncausal
\begin{column}{0.49\textwidth}
\begin{OrangeCard}[halign=center,top=10pt,bottom=10pt]
\CardTitle{Noncausal}

\vspace{0.04cm}

\begin{tikzpicture}[x=0.78cm,y=1.10cm,every node/.style={font=\scriptsize}]

  \draw[-{Latex[length=1.5mm]},thick,black]
    (-3.2,0) -- (3.4,0) node[right] {$t$};
  \draw[-{Latex[length=1.5mm]},thick,black]
    (0,-0.12) -- (0,1.35);

  % response before 0
  \draw[very thick,ThemeOrange,smooth]
    plot[domain=-2.8:0,samples=100] (\x,{exp(0.75*\x)});

  % zero after 0
  \draw[very thick,ThemeGray] (0,0) -- (2.8,0);

  \draw[densely dashed,ThemeGray] (0,0) -- (0,1.15);

  \node[below] at (0,-0.05) {$0$};
  \node[ThemeOrange,above] at (-1.45,0.58) {\(h(t)\)};
  \node[ThemeRed,font=\bfseries] at (3.25,0.68) {effect before cause};

\end{tikzpicture}

\vspace{0.06cm}

\(
h(t)\ne 0\qquad \text{for some }t<0.
\)

\end{OrangeCard}
\end{column}

\end{columns}

\end{frame}




% --------------------------------------------------------
\begin{frame}{Testing causality: the impulse response must be zero before \(0\)}

\begin{center}
\begin{GreenCard}[
  width=\textwidth,
  halign=center,
  top=12pt,
  bottom=12pt
]
\CardTitle{For continuous-time LTI systems}

\vspace{0.28cm}

{\Large
\(
\boxed{
\text{causal}
\quad\Longleftrightarrow\quad
h(t)=0\ \text{ for all }t<0
}
\)
}

% \vspace{0.12cm}

% \begin{OrangeBox}[
%   calloutshadedorange,
%   halign=center,
%   fontupper=\large
% ]{}
% Cannot respond before the impulse is applied.
% \end{OrangeBox}

\end{GreenCard}
\end{center}

\vspace{0.14cm}

\begin{center}
\begin{BlueCard}[
  width=0.88\textwidth,
  halign=center,
  top=10pt,
  bottom=10pt
]
\CardTitle{Discrete-time version}

\[
\boxed{
\text{causal}
\quad\Longleftrightarrow\quad
h[n]=0\ \text{ for all }n<0
}
\]

\end{BlueCard}
\end{center}

\end{frame}





% --------------------------------------------------------
\begin{frame}{Stability means bounded inputs stay bounded}

\begin{BlueBox}[
  calloutshadedblue,
  fontupper=\large,
  halign=center
]{}
A stable system should not turn a small, bounded input into an output that grows without limit.
\end{BlueBox}

\vspace{0.12cm}

\begin{columns}[c,totalwidth=\textwidth]

% stable idea
\begin{column}{0.49\textwidth}
\begin{GreenCard}[halign=center,top=10pt,bottom=10pt]
\CardTitle{Stable behavior}

\vspace{0.04cm}

\begin{tikzpicture}[x=0.75cm,y=1.05cm,every node/.style={font=\scriptsize}]

  \draw[-{Latex[length=1.5mm]},thick,black]
    (-2.2,0) -- (4.2,0) node[right] {$t$};
  \draw[-{Latex[length=1.5mm]},thick,black]
    (0,-0.10) -- (0,1.35);

  \draw[densely dashed,ThemeGray] (-2.0,1.0) -- (3.8,1.0);
  \draw[densely dashed,ThemeGray] (-2.0,-1.0) -- (3.8,-1.0);

  \draw[very thick,ThemeGreen,smooth]
    plot[domain=-2.0:3.8,samples=160] (\x,{0.75*sin(110*\x)});

  \node[ThemeGreen,above] at (2.3,0.65) {bounded output};
  \node[left] at (0,1) {\(B\)};
  \node[left] at (0,-1) {\(-B\)};

\end{tikzpicture}

\end{GreenCard}
\end{column}

% unstable idea
\begin{column}{0.49\textwidth}
\begin{OrangeCard}[halign=center,top=10pt,bottom=10pt]
\CardTitle{Unstable behavior}

\vspace{0.04cm}

\begin{tikzpicture}[x=0.75cm,y=0.54cm,every node/.style={font=\scriptsize}]

  \draw[-{Latex[length=1.5mm]},thick,black]
    (-2.2,0) -- (4.2,0) node[right] {$t$};
  \draw[-{Latex[length=1.5mm]},thick,black]
    (0,-0.20) -- (0,3.2);

  \draw[very thick,ThemeOrange,smooth]
    plot[domain=-1.2:2.55,samples=140] (\x,{0.30*exp(0.85*\x)});

  \node[ThemeRed,font=\bfseries] at (2.3,2.40) {keeps growing};

\end{tikzpicture}

\end{OrangeCard}
\end{column}

\end{columns}

\end{frame}



% --------------------------------------------------------
\begin{frame}{Stable LTI system: The total area under \(|h(t)|\) must be finite}

\begin{center}
\begin{GreenCard}[
  width=0.88\textwidth,
  halign=center,
  top=12pt,
  bottom=12pt
]
\CardTitle{Continuous-time LTI stability test}

\vspace{0.08cm}

{\Large
\(
\boxed{
\text{stable}
\quad\Longleftrightarrow\quad
\int_{-\infty}^{\infty}|h(t)|\,dt<\infty
}
\)
}

% \vspace{0.10cm}

% \begin{OrangeBox}[
%   calloutshadedorange,
%   halign=center,
%   fontupper=\large
% ]{}
% The total area under \(|h(t)|\) must be finite.
% \end{OrangeBox}

\end{GreenCard}
\end{center}

\vspace{0.12cm}

\begin{center}
\begin{BlueCard}[
  width=0.88\textwidth,
  halign=center,
  top=10pt,
  bottom=10pt
]
\CardTitle{Discrete-time version}

\[
\boxed{
\text{stable}
\quad\Longleftrightarrow\quad
\sum_{n=-\infty}^{\infty}|h[n]|<\infty
}
\]

\end{BlueCard}
\end{center}

\end{frame}






% --------------------------------------------------------
\begin{frame}{Practice: decide causality and stability}

\begin{OrangeBox}[
  calloutshadedorange,
  fontupper=\large,
  halign=center
]{}
For each impulse response, decide whether the LTI system is causal and stable.
\end{OrangeBox}

\vspace{0.14cm}

\begin{columns}[c,totalwidth=\textwidth]

\begin{column}{0.49\textwidth}
\begin{GreenCard}[halign=center,top=14pt,bottom=14pt]
\CardTitle{System 1}

\vspace{0.10cm}

{\Large
\[
h_1(t)=e^{2t}u(-2-t)
\]
}

% \vspace{0.10cm}

% {\normalsize
% Ask:
% \[
% h_1(t)=0\text{ for }t<0?
% \qquad
% \int |h_1(t)|dt<\infty?
% \]
% }

\end{GreenCard}
\end{column}

\begin{column}{0.49\textwidth}
\begin{BlueCard}[halign=center,top=14pt,bottom=14pt]
\CardTitle{System 2}

\vspace{0.10cm}

{\Large
\[
h_2(t)=t e^{-t}u(t)
\]
}

% \vspace{0.10cm}

% {\normalsize
% Ask:
% \[
% h_2(t)=0\text{ for }t<0?
% \qquad
% \int |h_2(t)|dt<\infty?
% \]
% }

\end{BlueCard}
\end{column}

\end{columns}

\end{frame}





% --------------------------------------------------------
\begin{frame}{Answer 1: \(h_1(t)=e^{2t}u(-2-t)\) is not causal}

\begin{center}
\begin{OrangeCard}[
  width=0.92\textwidth,
  halign=center,
  top=10pt,
  bottom=10pt
]
\CardTitle{Support of \(h_1(t)\)}

\vspace{0.06cm}

\[
u(-2-t)=1
\quad\Longleftrightarrow\quad
-2-t\ge 0
\quad\Longleftrightarrow\quad
t\le -2.
\]

\vspace{0.08cm}

\begin{tikzpicture}[x=0.88cm,y=1.20cm,every node/.style={font=\scriptsize}]

  \draw[-{Latex[length=1.5mm]},thick,black]
    (-5.0,0) -- (2.2,0) node[right] {$t$};

  \draw[-{Latex[length=1.5mm]},thick,black]
    (0,-0.10) -- (0,1.25);

  % zero after -2
  \draw[very thick,ThemeGray] (-2,0) -- (1.8,0);

  % nonzero before -2, vertically enlarged
  \draw[very thick,ThemeOrange,smooth]
    plot[domain=-4.7:-2,samples=120] (\x,{exp(2*(\x+2))});

  \draw[very thick,ThemeOrange] (-2,1) -- (-2,0);

  \node[ThemeGray,font=\large] at (-4.75,0.18) {$\cdots$};
  \node[ThemeGray,font=\large] at (1.65,0.18) {$\cdots$};

  \node[below] at (-2,-0.05) {$-2$};
  \node[below] at (0,-0.05) {$0$};

  \draw[densely dashed,ThemeRed,very thick] (0,0) -- (0,1.05);
  \node[ThemeRed,above] at (0,1.07) {impulse time};

  \node[ThemeOrange,above] at (-3.35,0.42) {\(t\le -2\)};
  \node[ThemeGray] at (-3.20,-0.32) {vertical scale enlarged};

\end{tikzpicture}

\vspace{0.08cm}

\begin{OrangeBox}[
  calloutshadedorange,
  halign=center,
  fontupper=\large
]{}
\(h_1(t)\) is nonzero before \(t=0\).  
So the system is \textbf{not causal}.
\end{OrangeBox}

\end{OrangeCard}
\end{center}

\end{frame}





% --------------------------------------------------------
\begin{frame}{Answer 1: \(h_1(t)=e^{2t}u(-2-t)\) is stable}

\begin{center}
\begin{GreenCard}[
  width=0.86\textwidth,
  halign=center,
  top=12pt,
  bottom=12pt
]
\CardTitle{Stability test}

\vspace{0.08cm}

\[
\int_{-\infty}^{\infty}|h_1(t)|\,dt
=
\int_{-\infty}^{-2} e^{2t}\,dt.
\]

\vspace{0.08cm}

\[
\int_{-\infty}^{-2} e^{2t}\,dt
=
\left[\frac{1}{2}e^{2t}\right]_{-\infty}^{-2}
=
\frac{1}{2}e^{-4}
<\infty.
\]

\vspace{0.12cm}

\begin{GreenBox}[
  calloutshadedgreen,
  halign=center,
  fontupper=\Large
]{}
Therefore, \(h_1(t)\) is \textbf{stable}.
\end{GreenBox}

\vspace{0.08cm}

\begin{BlueBox}[
  calloutshadedblue,
  halign=center,
  fontupper=\large
]{}
Final decision:
\[
h_1(t)=e^{2t}u(-2-t)
\quad\Rightarrow\quad
\textbf{not causal, stable}.
\]
\end{BlueBox}

\end{GreenCard}
\end{center}

\end{frame}




% --------------------------------------------------------
\begin{frame}{Answer 2: \(h_2(t)=t e^{-t}u(t)\) is causal}

\begin{center}
\begin{GreenCard}[
  width=0.92\textwidth,
  halign=center,
  top=10pt,
  bottom=10pt
]
\CardTitle{Support of \(h_2(t)\)}

\vspace{0.06cm}

\[
u(t)=0\quad\text{for }t<0.
\]

\vspace{0.08cm}

\begin{tikzpicture}[x=0.88cm,y=1.55cm,every node/.style={font=\scriptsize}]

  \draw[-{Latex[length=1.5mm]},thick,black]
    (-2.4,0) -- (5.2,0) node[right] {$t$};

  \draw[-{Latex[length=1.5mm]},thick,black]
    (0,-0.08) -- (0,1.10);

  % zero for t<0
  \draw[very thick,ThemeGreen] (-2.0,0) -- (0,0);

  % t e^{-t}, visually scaled
  \draw[very thick,ThemeGreen,smooth]
    plot[domain=0:4.7,samples=140] (\x,{2.55*\x*exp(-\x)});

  \node[ThemeGray,font=\large] at (-1.85,0.15) {$\cdots$};
  \node[ThemeGray,font=\large] at (4.75,0.15) {$\cdots$};

  \node[below] at (0,-0.04) {$0$};
  \node[ThemeGreen,above] at (2.0,0.80) {\(t e^{-t}u(t)\)};

  \draw[densely dashed,ThemeRed,very thick] (0,0) -- (0,0.95);
  \node[ThemeRed,above] at (0,0.97) {impulse time};

\end{tikzpicture}

\vspace{0.08cm}

\begin{GreenBox}[
  calloutshadedgreen,
  halign=center,
  fontupper=\large
]{}
\(h_2(t)\) is zero before \(t=0\).  
So the system is \textbf{causal}.
\end{GreenBox}

\end{GreenCard}
\end{center}

\end{frame}







% --------------------------------------------------------
\begin{frame}{Answer 2: \(h_2(t)=t e^{-t}u(t)\) is stable}

\begin{center}
\begin{GreenCard}[
  width=0.86\textwidth,
  halign=center,
  top=12pt,
  bottom=12pt
]
\CardTitle{Stability test}

\vspace{0.08cm}

\[
\int_{-\infty}^{\infty}|h_2(t)|\,dt
=
\int_{0}^{\infty}t e^{-t}\,dt.
\]

\vspace{0.08cm}

\[
\int_{0}^{\infty}t e^{-t}\,dt
=
1
<\infty.
\]

\vspace{0.12cm}

\begin{GreenBox}[
  calloutshadedgreen,
  halign=center,
  fontupper=\Large
]{}
Therefore, \(h_2(t)\) is \textbf{stable}.
\end{GreenBox}

\vspace{0.08cm}

\begin{BlueBox}[
  calloutshadedblue,
  halign=center,
  fontupper=\large
]{}
Final decision:
\[
h_2(t)=t e^{-t}u(t)
\quad\Rightarrow\quad
\textbf{causal, stable}.
\]
\end{BlueBox}

\end{GreenCard}
\end{center}

\end{frame}








\end{document}